\chapter{OFDM and Multicarrier Transmission}
\label{ch:ofdm}

\begin{chapterintro}
Orthogonal frequency-division multiplexing is the dominant waveform of the early
twenty-first century. It carries Wi-Fi from 802.11a to Wi-Fi~7, the LTE and 5G NR
downlinks, digital television and radio on four continents, ADSL and VDSL over telephone
wires, DOCSIS~3.1 over cable, HomePlug over the power line, and a growing share of optical and
visible-light links. The idea is to split a wide channel into hundreds or thousands of narrow
subcarriers, each so narrow that the channel looks flat across it, and to generate and
separate them all with one FFT. A cyclic prefix turns the channel's convolution into a
multiplication, so equalisation becomes one complex multiply per subcarrier. This chapter
starts with the history, from Chang's 1966 Bell Labs paper to LTE's decision in 2005. It
derives OFDM from the DFT, proves the cyclic-prefix property in matrix form, and shows how
subcarrier spacing, prefix length and FFT size follow from the delay spread and the
Doppler. The worked designs are 802.11a, LTE, the 5G NR numerologies and the DVB-T2 32k mode.
The middle of the chapter covers the resource grid, pilots and channel estimation, and then
what breaks OFDM: timing and frequency offsets, Doppler, phase noise, sampling-clock error,
the peak-to-average power ratio and out-of-band emission. The chapter then covers coded OFDM
and frequency diversity, bit loading and water-filling in DSL, OFDMA and MIMO-OFDM,
single-frequency networks, OFDM in Wi-Fi~6 and~7, on power lines and in optics, and closes
with a first look at OFDM radar.
\end{chapterintro}

%======================================================================
\section{Why multicarrier?}
\label{sec:ch17:why}
%======================================================================

\subsection{The problem: a wideband channel}

Chapter~\ref{ch:channels} described the mobile radio channel as a sum of echoes. In a city the
echoes arrive over a few microseconds; in hilly terrain over ten or twenty. Chapter~\ref{ch:equalization}
then showed what this does to a single-carrier signal. At 20 million symbols per second, each
symbol lasts 50\,ns. A 5\,$\mu$s delay spread therefore smears each symbol over about a hundred of its
neighbours, and a 10\,$\mu$s spread over two hundred. A time-domain equaliser for such a channel needs
hundreds of taps, adapts slowly because its error surface has a large eigenvalue spread, and costs
hundreds of complex multiplies per symbol. Maximum-likelihood sequence detection is out of the
question. This is the \emph{complexity wall} of Section~\ref{sec:ch12:fde}.

The multicarrier answer is to divide and conquer. Split the 20\,MHz band into $N$ subchannels.
With $N=1024$, each subchannel is about 20\,kHz wide and carries symbols about 50\,$\mu$s long. A
5\,$\mu$s echo now smears each symbol by a tenth of its duration instead of by a hundred symbols, and
across a 20\,kHz slice of spectrum the channel's frequency response is almost constant. Each
subchannel sees a \emph{flat} channel, a single complex gain, and needs only a one-tap equaliser.
The ISI problem has not disappeared. It has been made small enough that a short guard interval
can absorb it completely.

\subsection{From guard bands to overlapping subcarriers}

Frequency-division multiplexing (FDM) is as old as radio. The classical way to put several data
signals side by side is to band-limit each one with a filter and to leave a guard band between
neighbours, so that a realisable receive filter can pick one out
(Figure~\ref{fig:ch17_fdm_vs_ofdm}b). The telephone network's analog carrier systems stacked voice
channels this way, 4\,kHz apart, for most of the twentieth century
(Chapter~\ref{ch:history}). Applied to data, FDM has two drawbacks. The guard bands waste spectrum,
typically 20--50\%. And a bank of hundreds of sharp analog filters, each with its own oscillator,
is expensive and hard to keep aligned.

The key observation is that the subchannels do not have to be spectrally disjoint to be separable.
They have to be \emph{orthogonal}. Consider $N$ complex exponentials on a symbol interval of length $T$,
\begin{equation}
\phi_k(t)=\frac{1}{\sqrt T}\,e^{j2\pi k\Delta f\,t},\qquad 0\le t<T,\qquad k=0,1,\dots,N-1,
\label{eq:ch17:subcarrier}
\end{equation}
with spacing $\Delta f=1/T$. Each one completes a whole number of cycles in $T$
(Figure~\ref{fig:ch17_fdm_vs_ofdm}a), and their inner products are
\begin{equation}
\langle\phi_k,\phi_l\rangle=\frac1T\int_0^T e^{j2\pi(k-l)t/T}\,dt=
\begin{cases}1,&k=l,\\[2pt] \dfrac{e^{j2\pi(k-l)}-1}{j2\pi(k-l)}=0,&k\ne l.\end{cases}
\label{eq:ch17:orth}
\end{equation}
The subcarriers are orthonormal. A correlator matched to subcarrier $k$ therefore sees none of the
others, even though their spectra overlap heavily. In the frequency domain, each time-limited
subcarrier has a spectrum $T\sinc\big((f-k\Delta f)T\big)$ whose zeros fall at every \emph{other}
subcarrier's centre frequency (Figure~\ref{fig:ch17_fdm_vs_ofdm}c). Sampled at $f=k\Delta f$, only
subcarrier $k$ contributes. This is \term{orthogonal frequency-division multiplexing} (OFDM).

\mdcfig{ch17_fdm_vs_ofdm}{From FDM to OFDM. (a) Subcarriers spaced by $1/T$ complete a whole number of
cycles in the symbol interval $T$, which makes them orthogonal. (b) Classical FDM separates band-limited
channels with guard bands and needs one sharp filter per channel. (c) OFDM overlaps the $\sinc$-shaped
subcarrier spectra; at each subcarrier's centre frequency (dots) all the others pass through zero.}

The spectral efficiency gain is easy to count. $N$ subcarriers at spacing $1/T$ occupy a total
bandwidth of about $(N+1)/T$ (the main lobes of the two edge subcarriers stick out by $1/T$ in
total) while carrying $N$ complex symbols every $T$ seconds. The efficiency approaches one complex
symbol per second per hertz, the Nyquist limit of Chapter~\ref{ch:baseband}, without needing a
single sharp filter. OFDM is to FDM what the Nyquist pulse is to a rectangular spectrum: it
replaces a brick wall with a carefully arranged overlap that cancels exactly at the sampling
points. The duality is exact: a Nyquist pulse is zero at the other symbols' sampling times; an
OFDM subcarrier's spectrum is zero at the other subcarriers' frequencies.

\begin{keyidea}[OFDM turns one hard channel into many easy ones]
A frequency-selective channel with a long impulse response is hard to equalise in the time domain.
OFDM transmits $N$ narrowband signals in parallel, each seeing an essentially flat channel. With a
cyclic prefix (Section~\ref{sec:ch17:cp}) the conversion is exact: the channel becomes $N$
independent scalar channels $Y_k=H_kX_k+W_k$. Almost everything else in this chapter is either a
consequence of this identity or a cost paid to keep it true.
\end{keyidea}

\begin{history}[From Kineplex to Chang's orthogonal channels]
Parallel-tone data transmission predates the theory. Collins Radio's \emph{Kineplex} system of the
late 1950s sent data over HF radio on 20 tones, each carrying differentially encoded phase
shifts at a low symbol rate so that skywave multipath echoes of a few milliseconds were short
compared with a symbol. The US military's KATHRYN HF modem of the mid-1960s (described by
Zimmerman and Kirsch in 1967) went further, generating its 34 subchannels with what was
effectively a discrete Fourier transform implemented in analog hardware. The theory came from
Robert W.~Chang at Bell Labs, whose 1966 \emph{Bell System Technical Journal} paper ``Synthesis of
band-limited orthogonal signals for multichannel data transmission'' gave the conditions under
which overlapping, band-limited subchannels could be transmitted with neither intersymbol nor
inter-channel interference. Chang patented the scheme (US patent 3,488,445, granted 1970). In
1967 Burton Saltzberg, also at Bell Labs, analysed a version with staggered (offset) QAM on each
subcarrier that tolerates tighter spectral shaping; his idea returned forty years later as
FBMC/OQAM (Section~\ref{sec:ch17:oob}). Both schemes needed a bank of analog oscillators and
filters, and for more than a decade they remained a curiosity.
\end{history}

%======================================================================
\section{OFDM with the discrete Fourier transform}
\label{sec:ch17:dft}
%======================================================================

\subsection{The transmitter is an inverse DFT}

The complex baseband OFDM symbol carrying data symbols $X_0,\dots,X_{N-1}$ (drawn from QAM
constellations, one per subcarrier) is
\begin{equation}
s(t)=\sum_{k=0}^{N-1}X_k\,e^{j2\pi k\Delta f\,t},\qquad 0\le t<T=\frac1{\Delta f}.
\label{eq:ch17:ofdmct}
\end{equation}
Its bandwidth is about $N\Delta f$, so by the sampling theorem it is represented by $N$ samples
taken at $t=nT/N$, $n=0,\dots,N-1$:
\begin{equation}
x[n]=s\!\left(\frac{nT}{N}\right)=\sum_{k=0}^{N-1}X_k\,e^{j2\pi kn/N}=N\cdot\mathrm{IDFT}\{X_k\}[n].
\label{eq:ch17:idft}
\end{equation}
This is the observation of Stephen Weinstein and Paul Ebert (Bell Labs, 1971): the whole bank of
modulators is an inverse discrete Fourier transform, and the bank of receive correlators is a
forward DFT,
\begin{equation}
\frac1N\sum_{n=0}^{N-1}x[n]\,e^{-j2\pi kn/N}=X_k.
\label{eq:ch17:dftrx}
\end{equation}
The orthogonality of~\eqref{eq:ch17:orth} is now the orthogonality of the DFT basis vectors, which
holds exactly in discrete time. With the FFT, generating $N$ subcarriers costs about
$\tfrac N2\log_2N$ complex multiplies per OFDM symbol, or $\tfrac12\log_2N$ per data symbol:
5 for $N=1024$. No analog filter bank is needed.

In what follows we use the unitary normalisation $\vect x=\vect F^H\vect X$ and $\vect X=\vect F\vect x$,
where $\vect F$ is the $N\times N$ DFT matrix with entries $[\vect F]_{kn}=e^{-j2\pi kn/N}/\sqrt N$.
It preserves energy (Parseval), so the average power per sample equals the average power per
subcarrier. The \texttt{commlib.ofdm} functions used by the labs follow this convention.

\begin{history}[Weinstein and Ebert, Peled and Ruiz, and Cimini]
In 1971 Stephen Weinstein and Paul Ebert of Bell Labs published ``Data transmission by
frequency-division multiplexing using the discrete Fourier transform'' in the \emph{IEEE Transactions
on Communication Technology}. The FFT algorithm of Cooley and Tukey was six years old, and they saw
that it replaced the entire oscillator bank. They also proposed a guard time between symbols to
absorb channel dispersion, although with an empty guard, which loses orthogonality on a dispersive
channel. In 1980 Abraham Peled and Antonio Ruiz of IBM filled the guard with a \emph{cyclic extension}
of the symbol, the cyclic prefix used by every OFDM system since. In 1985 Leonard Cimini of Bell Labs
analysed OFDM for the mobile radio channel, with pilot-based channel estimation and a one-tap
equaliser, showing it could handle the frequency-selective fading that was defeating equalisers of
the time. The name ``OFDM'' became standard in the 1980s; John Bingham's 1990 \emph{IEEE Communications
Magazine} article ``Multicarrier modulation for data transmission: an idea whose time has come''
announced that cheap digital signal processing had finally made it practical. Weinstein later wrote
a short history of the whole development for the same magazine (2009), listed in the further reading.
\end{history}

\subsection{Guard subcarriers, the DC null and oversampling}

A practical OFDM system never uses all $N$ bins. The bins near $\pm N/2$ (the edges of the sampled
band) are left empty as \term{guard subcarriers}. They give the DAC's reconstruction filter and the
receiver's anti-aliasing filter a transition band, and they keep the sidelobes of the occupied
subcarriers away from the neighbouring channel. The \term{DC subcarrier} is also usually left
empty, because direct-conversion receivers have a DC offset and local-oscillator leakage that lands
exactly there (Chapter~\ref{ch:sdr}). An 802.11a transmitter uses 52 of its 64 bins: subcarriers
$-26$ to $+26$ without 0, occupying 16.6\,MHz of a 20\,MHz channel
(Figure~\ref{fig:ch17_ofdm_symbol}b). An LTE carrier of 20\,MHz uses 1200 of 2048 bins, 18\,MHz of 20.

Leaving the outer bins empty has a second benefit: the IFFT output is then an \emph{oversampled}
version of the OFDM signal. Zero-padding the frequency-domain vector to $LN$ bins produces $L$
times as many time samples of the same waveform. This is how the figures in this chapter were
drawn, and it matters for the peak-to-average ratio (Section~\ref{sec:ch17:papr}), whose true peaks
fall between Nyquist-rate samples.

\mdcfig{ch17_ofdm_symbol}{An 802.11a-like OFDM signal (64-point FFT at 20\,MS/s, 52 QPSK or 16-QAM
subcarriers). (a) The magnitude of three consecutive 4\,$\mu$s symbols. Each begins with a 0.8\,$\mu$s cyclic
prefix (red) that repeats the symbol's last 0.8\,$\mu$s (green). The waveform looks like noise, with
occasional high peaks. (b) The power spectral density: flat across the occupied band, with the empty DC
subcarrier visible as a notch. Outside the band the $\sinc$ sidelobes of the rectangular symbols decay
only as $1/f^2$.}

\subsection{The complete transceiver}

Figure~\ref{fig:ch17_chain} shows a complete OFDM link. The transmitter encodes and interleaves the
bits (Section~\ref{sec:ch17:coded}), maps them onto QAM symbols, and places the symbols on the
data subcarriers of a frequency-domain grid, together with known \emph{pilot} symbols
(Section~\ref{sec:ch17:grid}). The IFFT converts each column of the grid into $N$ time samples. A
cyclic prefix is prepended, the symbols are concatenated, optionally windowed or filtered to
control the spectrum (Section~\ref{sec:ch17:oob}), converted to analog and up-converted. The
receiver reverses the chain. It first finds the symbol timing and the carrier frequency offset
(Section~\ref{sec:ch17:sync} and Chapter~\ref{ch:sync}), then strips the prefix, takes an FFT,
estimates the channel from the pilots, equalises each subcarrier with one complex multiply, and
passes soft bit decisions to the decoder.

\begin{figure}[htbp]\centering
\begin{tikzpicture}[node distance=0.34cm and 0.34cm,
  blk/.style={draw=navy,fill=navy!6,thick,minimum height=0.9cm,minimum width=1.2cm,align=center,font=\sffamily\scriptsize},
  key/.style={blk,fill=accent!8,draw=accent},
  arr/.style={-{Stealth[length=1.8mm]},thick,navy}]
% transmitter
\node[font=\sffamily\scriptsize\bfseries,text=navy,anchor=west] at (-0.7,1.05) {TRANSMITTER};
\node[blk] (enc) {FEC and\\interleave};
\node[blk,right=of enc] (map) {QAM\\map};
\node[blk,right=of map] (grid) {Map to\\subcarriers\\+ pilots};
\node[key,right=of grid] (ifft) {$N$-point\\IFFT};
\node[key,right=of ifft] (cp) {Add\\cyclic\\prefix};
\node[blk,right=of cp] (win) {Window\\or filter};
\node[blk,right=of win] (dac) {DAC,\\up-convert,\\PA};
\draw[arr] (enc)--(map); \draw[arr] (map)--(grid); \draw[arr] (grid)--(ifft);
\draw[arr] (ifft)--(cp); \draw[arr] (cp)--(win); \draw[arr] (win)--(dac);
\node[left=0.25cm of enc,font=\sffamily\scriptsize] (bits) {bits};
\draw[arr] (bits)--(enc);
% channel
\node[blk,below=0.55cm of dac,fill=accent!4,draw=accent,minimum width=1.2cm] (ch) {Channel\\$h[n]$\\+ noise};
\draw[arr,accent] (dac)--(ch);
% receiver (right to left)
\node[blk,below=0.55cm of ch] (adc) {Down-\\convert,\\ADC};
\node[blk,left=of adc] (sync) {Timing\\and CFO\\sync};
\node[key,left=of sync] (rcp) {Remove\\cyclic\\prefix};
\node[key,left=of rcp] (fft) {$N$-point\\FFT};
\node[key,left=of fft] (eq) {One-tap\\equaliser\\$\hat X_k=Y_k/\hat H_k$};
\node[blk,left=of eq] (demap) {Soft\\demap\\(LLRs)};
\node[blk,left=of demap] (dec) {De-inter-\\leave and\\decode};
\node[blk,above=0.28cm of eq,fill=practbg,draw=practc,minimum height=0.6cm] (est) {Channel estimation (pilots)};
\draw[arr,accent] (ch)--(adc);
\draw[arr] (adc)--(sync); \draw[arr] (sync)--(rcp); \draw[arr] (rcp)--(fft);
\draw[arr] (fft)--(eq); \draw[arr] (eq)--(demap); \draw[arr] (demap)--(dec);
\draw[arr,practc] (fft.north) |- (est.east);
\draw[arr,practc] (est.south)--(eq.north);
\node[font=\sffamily\scriptsize\bfseries,text=navy,anchor=west] at ([yshift=-0.75cm]dec.south west) {RECEIVER};
\node[below=0.25cm of dec,font=\sffamily\scriptsize] (bout) {bits};
\draw[arr] (dec)--(bout);
\end{tikzpicture}
\caption{An OFDM transceiver. The red blocks are the ones specific to OFDM: the IFFT, the cyclic prefix,
the FFT and the per-subcarrier equaliser. Everything else---coding, mapping, synchronisation,
channel estimation, RF---has a single-carrier counterpart, but in OFDM it all operates on a
time--frequency grid of resource elements.}
\label{fig:ch17_chain}
\end{figure}


%======================================================================
\section{The cyclic prefix}
\label{sec:ch17:cp}
%======================================================================

\subsection{What goes wrong without it}

Pass one OFDM symbol through a channel with impulse response $h[0],\dots,h[\nu]$ (memory $\nu$
samples). The receiver sees the linear convolution $y[n]=\sum_mh[m]x[n-m]+w[n]$, which is $\nu$
samples longer than the symbol. Two things go wrong. The tail of each symbol spills into the next,
which is intersymbol interference in the familiar sense. More subtly, within one symbol the
subcarriers stop being orthogonal. A delayed copy of subcarrier $k$ starts part-way through the FFT
window, so within the window it is a truncated sinusoid rather than a whole number of cycles, and
its spectrum leaks into every other bin. This is \term{inter-carrier interference} (ICI).
Weinstein and Ebert's empty guard interval removes the first problem but not the second.

\subsection{Linear convolution made circular}

The cure is to copy the last $N_{\text{cp}}$ samples of each symbol and transmit them first. The
transmitted block is
\begin{equation}
\tilde{\vect x}=\big[x[N-N_{\text{cp}}],\dots,x[N-1],\;x[0],x[1],\dots,x[N-1]\big]^T,
\end{equation}
of length $N+N_{\text{cp}}$. Suppose $N_{\text{cp}}\ge\nu$. The receiver discards the first
$N_{\text{cp}}$ received samples of each block and keeps the next $N$. Every kept sample
$y[n]$, $0\le n\le N-1$, is a combination of $x[n],x[n-1],\dots,x[n-\nu]$, and whenever an index
$n-m$ is negative the sample that arrives is the prefix copy $x[n-m+N]$. So
\begin{equation}
y[n]=\sum_{m=0}^{\nu}h[m]\,x\big[(n-m)\bmod N\big]+w[n],\qquad n=0,\dots,N-1,
\label{eq:ch17:circ}
\end{equation}
which is the \emph{circular} convolution of $\vect x$ with $\vect h$. The previous symbol's
samples, which are the source of ISI, fall only on the discarded prefix
(Figure~\ref{fig:ch17_cp_tikz}). And circular convolution is exactly what the DFT turns into
multiplication (Chapter~\ref{ch:signals}):
\begin{equation}
\boxed{\;Y_k=H_k\,X_k+W_k,\qquad H_k=\sum_{m=0}^{\nu}h[m]\,e^{-j2\pi km/N},\qquad k=0,\dots,N-1.\;}
\label{eq:ch17:onetap}
\end{equation}
Each subcarrier sees a single complex gain $H_k$, the channel's frequency response sampled at that
subcarrier, and additive noise $W_k$ with the same variance as $w[n]$ because $\vect F$ is unitary.
A frequency-selective channel has become $N$ parallel flat-fading channels with no interference
between them.

\begin{figure}[htbp]\centering
\begin{tikzpicture}[x=0.86cm,y=0.62cm,font=\sffamily\scriptsize]
% direct path
\node[anchor=east] at (-0.2,2.5) {direct path};
\node[anchor=east] at (-0.2,1.3) {echo, delay $\tau$};
\foreach \s/\c in {0/navy,1/accent}{
  \pgfmathsetmacro{\xa}{\s*5.4}
  \fill[\c!30] (\xa,2.1) rectangle ++(1.1,0.8);
  \fill[\c!10] (\xa+1.1,2.1) rectangle ++(4.3,0.8);
  \draw[\c,thick] (\xa,2.1) rectangle ++(5.4,0.8);
  \draw[\c] (\xa+1.1,2.1) -- ++(0,0.8);
  \node at (\xa+0.55,2.5) {CP};
  \node at (\xa+3.25,2.5) {symbol \s};
  \fill[\c!30] (\xa+0.7,0.9) rectangle ++(1.1,0.8);
  \fill[\c!10] (\xa+1.8,0.9) rectangle ++(4.3,0.8);
  \draw[\c,thick,dashed] (\xa+0.7,0.9) rectangle ++(5.4,0.8);
}
\node at (6.95,1.3) {\textcolor{accent}{CP}};
\node at (9.35,1.3) {\textcolor{accent}{symbol 1 (echo)}};
% ISI region
\fill[pattern=north east lines,pattern color=accent] (5.4,0.9) rectangle (6.1,1.7);
\draw[decorate,decoration={brace,amplitude=4pt,mirror}] (5.4,0.8) -- (6.1,0.8) coordinate[midway,below=4pt] (isi); \node[anchor=north east,text=accent] at ([xshift=0.35cm]isi) {ISI from symbol 0};
% FFT window
\draw[practc,very thick] (6.5,-0.4) rectangle (10.8,3.3);
\node[practc,anchor=south] at (8.65,3.3) {FFT window ($N$ samples)};
\draw[navy,thick,{Stealth}-{Stealth}] (5.4,3.55) -- (6.5,3.55) node[midway,above]{$N_{\text{cp}}$};
\draw[accent,{Stealth}-{Stealth}] (5.4,1.9) -- (6.1,1.9) node[midway,left=6pt,inner sep=1pt]{$\tau$};
\node[align=left,anchor=west,text width=2.8cm] at (12.2,1.8) {Inside the window every\\path contributes whole\\cycles of symbol~1 only:\\ISI and ICI vanish if\\$\tau\le N_{\text{cp}}$.};
\end{tikzpicture}
\caption{Why the cyclic prefix works. An echo delayed by $\tau$ lets the end of symbol~0 overlap the
start of symbol~1 (hatched), but only inside the prefix, which the receiver discards. Within the FFT
window, the echo consists of a cyclically shifted copy of symbol~1, so each subcarrier still has a whole
number of cycles and orthogonality is preserved. The echo appears in the subcarrier domain only as a phase
rotation $e^{-j2\pi k\Delta f\tau}$, absorbed into $H_k$.}
\label{fig:ch17_cp_tikz}
\end{figure}

\subsection{The matrix view: circulant channels are diagonalised by the DFT}

The same result has an elegant linear-algebra form that generalises to MIMO and to the analysis of
every impairment later in the chapter. With the prefix inserted and removed, the channel from
$\vect x$ to $\vect y$ is the $N\times N$ \term{circulant matrix}
\begin{equation}
\vect H_c=\begin{bmatrix}
h[0] & 0 & \cdots & 0 & h[\nu] & \cdots & h[1]\\
h[1] & h[0] & \cdots & 0 & 0 & \ddots & \vdots\\
\vdots & & \ddots & & & \ddots & h[\nu]\\
h[\nu] & \cdots & & h[0] & & & 0\\
\vdots & \ddots & & & \ddots & & \vdots\\
0 & \cdots & h[\nu] & \cdots & & h[1] & h[0]
\end{bmatrix},
\qquad \vect y=\vect H_c\vect x+\vect w,
\end{equation}
each of whose rows is the previous row shifted cyclically one place to the right. Every circulant
matrix has the same eigenvectors, the columns of $\vect F^H$, whatever its entries. To see why, apply
$\vect H_c$ to the $k$th Fourier vector $\vect f_k=[1,e^{j2\pi k/N},\dots,e^{j2\pi k(N-1)/N}]^T/\sqrt N$:
row $n$ gives $\sum_mh[m]e^{j2\pi k(n-m)/N}/\sqrt N=H_k\,[\vect f_k]_n$. The eigenvalues are the
DFT of the first column. Hence
\begin{equation}
\vect H_c=\vect F^H\bm{\Lambda}\vect F,\qquad \bm{\Lambda}=\operatorname{diag}(H_0,\dots,H_{N-1}).
\label{eq:ch17:eig}
\end{equation}
The transmitter sends $\vect x=\vect F^H\vect X$, the receiver computes
\begin{equation}
\vect Y=\vect F\vect y=\vect F\vect F^H\bm{\Lambda}\vect F\vect F^H\vect X+\vect F\vect w=\bm{\Lambda}\vect X+\vect W,
\end{equation}
which is~\eqref{eq:ch17:onetap} in matrix form. OFDM is simply the eigen-decomposition of the
channel: it transmits along the channel's eigenvectors, which pass through unchanged apart from a
scaling by the eigenvalues. Seen this way, OFDM is the practical realisation of the
\emph{parallel Gaussian channels} of information theory (Chapter~\ref{ch:infotheory}), and the
eigenvectors do not depend on the channel, so the transmitter does not need to know it. Without the
prefix, $\vect H$ would be Toeplitz rather than circulant, and the DFT would diagonalise it only
approximately, leaving off-diagonal terms: the ICI.

\begin{keyidea}[The cyclic prefix buys channel-independent eigenvectors]
Any linear channel could be diagonalised by its own singular value decomposition, but the transmitter
would have to know the channel and recompute the transform for every channel realisation. The cyclic
prefix makes the channel matrix circulant, and all circulant matrices share the DFT as their
eigenbasis. One fixed, fast transform therefore diagonalises every channel whose memory is shorter
than the prefix.
\end{keyidea}

\subsection{One-tap equalisation}

Once~\eqref{eq:ch17:onetap} holds, equalisation is a scalar problem on each subcarrier. The
zero-forcing equaliser is $\hat X_k=Y_k/H_k$; the MMSE equaliser is
\begin{equation}
\hat X_k=\frac{H_k^*}{|H_k|^2+N_0/E_s}\,Y_k,
\end{equation}
the scalar version of the Wiener solution of Chapter~\ref{ch:equalization}. The two differ only on deeply
faded subcarriers, and for a coded system the choice hardly matters, because what the decoder
actually needs is not $\hat X_k$ but the per-bit log-likelihood ratios computed from $Y_k$, $H_k$ and
$N_0$ (Chapter~\ref{ch:modulation}). The effective SNR on subcarrier $k$ is
$|H_k|^2E_s/N_0$, so the LLRs of bits on a faded subcarrier are automatically small, and the decoder
treats them as the unreliable guesses they are.

Figure~\ref{fig:ch17_one_tap} shows the result on the LTE EVA channel with 600 subcarriers of 15\,kHz.
The channel varies by more than 25\,dB across the 9\,MHz band. The received subcarrier values are
scattered in amplitude and phase, and one complex division per subcarrier restores a clean 16-QAM
constellation. The residual scatter comes from the few subcarriers near the deep notches, where the
equaliser divides noise by a small number. Uncoded, those subcarriers dominate the error rate;
Section~\ref{sec:ch17:coded} shows how coding across subcarriers removes the problem.

\mdcfig{ch17_one_tap}{One-tap equalisation on a static EVA channel (Lab~7 setting: 1024-point FFT at
15.36\,MS/s, 600 subcarriers, 72-sample prefix, SNR 30\,dB). (a) The channel's frequency response
across the occupied band. (b) Received 16-QAM subcarrier values: each has been scaled and rotated by its
own $H_k$. (c) After dividing by $H_k$, the constellation is restored; the scattered points come from
subcarriers near the notches, where noise is amplified.}

\subsection{When the prefix is too short}

Real channels do not stop abruptly at $\nu$ samples; their power delay profiles have tails. If the
channel extends beyond the prefix, the energy in the excess taps causes both ISI (from the previous
symbol) and ICI (from the loss of circularity). A path of relative power $p_i$ arriving $\tau_i$ after
the first contributes, if $\tau_i>T_{\text{cp}}$, a fraction $(\tau_i-T_{\text{cp}})/T$ of its
energy as ISI from the previous symbol, and roughly the same again as ICI, because that same fraction
of the current symbol is missing from its circular convolution. For small excesses, the
signal-to-interference ratio is therefore approximately
\begin{equation}
\frac{S}{I}\approx\left[\,2\sum_i p_i\,\frac{(\tau_i-T_{\text{cp}})^+}{T}\right]^{-1},
\label{eq:ch17:cpsir}
\end{equation}
with $(u)^+=\max(u,0)$ and $\sum_ip_i=1$. Figure~\ref{fig:ch17_cp_length} measures the effect for the
three LTE profiles at an SNR of 45\,dB. For each profile, the interference falls to the noise floor
only when the prefix reaches the maximum excess delay; with a shorter prefix, the error floor is set
by the late paths. With no prefix at all, \eqref{eq:ch17:cpsir} gives $T/(2\bar\tau)$, where $\bar\tau$ is
the mean excess delay: for ETU ($\bar\tau\approx0.56\,\mu$s, $T=66.7\,\mu$s) about 18\,dB, as
measured. The LTE normal prefix of 4.69\,$\mu$s covers EPA and EVA with a large margin and ETU almost
completely. Only ETU's last tap, at 5\,$\mu$s with about 3\% of the power, sticks out, by 0.31\,$\mu$s,
and~\eqref{eq:ch17:cpsir} gives $S/I\approx1/(2\times0.031\times0.0047)$, about 35\,dB, again as measured.

\mdcfig[0.75]{ch17_cp_length}{Interference plus noise on the subcarriers as a function of prefix length,
for the three LTE channel models (static realisations, averaged; 1024-point FFT at 15.36\,MS/s; SNR
45\,dB). Dotted lines mark each model's maximum excess delay. Once the prefix covers it, only noise
remains; with a shorter prefix, the late echoes set an error floor no transmit power can overcome.}

\begin{worked}[The cost of the prefix in 802.11a and LTE]
The prefix is pure overhead in both time and energy. The energy spent on it is discarded by the
receiver, so the useful $E_s/N_0$ is lower than the transmitted one by
$10\log_{10}\big((N+N_{\text{cp}})/N\big)$\,dB, and the symbol rate is reduced by the same ratio.

\emph{802.11a:} $N=64$, $N_{\text{cp}}=16$ at 20\,MS/s ($T_{\text{cp}}=0.8\,\mu$s,
$T=3.2\,\mu$s). The ratio is $80/64=1.25$: 20\% of the air time and 0.97\,dB of SNR.

\emph{LTE, normal prefix:} $N=2048$ at 30.72\,MS/s ($T=66.7\,\mu$s), with seven symbols per 0.5\,ms
slot, the first having a prefix of 160 samples and the other six 144. The slot holds
$7\times2048=14336$ useful samples out of 15360, an efficiency of 93.3\%: 6.7\% overhead and
0.30\,dB.

\emph{LTE, extended prefix:} six symbols per slot, each with a 512-sample (16.7\,$\mu$s) prefix:
$6\times2048=12288$ of 15360 samples, 20\% overhead, 0.97\,dB.

The pattern is a direct trade: a prefix long enough for large cells or broadcast costs the same 20\%
as Wi-Fi's short prefix at a subcarrier spacing twenty times wider. Section~\ref{sec:ch17:design}
turns this trade into a design procedure.
\end{worked}

\subsection{Alternatives to the cyclic prefix}

The prefix is not the only guard that makes the channel circulant. \term{Zero-padded OFDM} (ZP-OFDM)
transmits nothing in the guard interval and, at the receiver, adds the $\nu$ samples that spill past
the end of the FFT window back onto its beginning (overlap-add), which reproduces the circular
convolution. It saves the transmit power spent on the prefix and allows the receiver to recover
symbols even on subcarriers where $H_k=0$, at the cost of slightly more noise; the WiMedia multiband
OFDM ultra-wideband standard (ECMA-368) used it. A \term{unique word} or \term{known-symbol padding}
fills the guard with a fixed sequence that can double as a pilot; the single-carrier PHYs of IEEE
802.11ad/ay use Golay sequences this way. \term{Single-carrier frequency-domain equalisation}
(SC-FDE, Section~\ref{sec:ch12:fde}) keeps the prefix but moves the IFFT from the transmitter to the
receiver, and the uplinks of LTE and NR use a hybrid of the two, DFT-spread OFDM
(Section~\ref{sec:ch17:dfts}).

%======================================================================
\section{Designing an OFDM system}
\label{sec:ch17:design}
%======================================================================

\subsection{Four numbers and three constraints}

An OFDM physical layer is defined, to first order, by four numbers: the occupied bandwidth $B$, the
subcarrier spacing $\Delta f=1/T$, the prefix duration $T_{\text{cp}}$, and the FFT size $N$ (which,
with $\Delta f$, fixes the sample rate $f_s=N\Delta f$). Three physical constraints pull them in
different directions.

\paragraph{The prefix must cover the delay spread.} $T_{\text{cp}}$ must exceed the channel's maximum
excess delay, not its RMS delay spread (the pitfall at the end of Section~\ref{sec:ch11:standards}), plus
the timing uncertainty of the receiver and the delay of the transmit and receive filters. Indoor
channels need a few hundred nanoseconds; urban macro-cells 3--5\,$\mu$s; hilly terrain and
single-frequency broadcast networks tens to hundreds of microseconds.

\paragraph{The overhead must be tolerable.} The prefix wastes a fraction
$T_{\text{cp}}/(T+T_{\text{cp}})$ of the time and energy. To keep this below 10--20\%, the useful
symbol must be five to ten times longer than the prefix, and so $\Delta f=1/T$ must be correspondingly
small.

\paragraph{The spacing must dwarf the Doppler and the oscillator impairments.} A channel that changes
during a symbol destroys orthogonality. Section~\ref{sec:ch17:sync} shows that a Doppler spread with
maximum frequency $f_D$ causes ICI of relative power about $(\pi f_DT)^2/6$; residual carrier offset
and phase noise add more. To keep ICI 25--30\,dB below the signal, $f_DT$ must be below about
0.03--0.05, which pushes $\Delta f$ \emph{up}. Wider spacing also shortens the symbol, which lowers the
latency and makes the system more robust to phase noise, a dominant concern at millimetre-wave
frequencies.

The FFT size then follows from $N\ge B/\Delta f$, rounded up to a convenient power of two with room
for guard subcarriers, and the sample rate from $N\Delta f$. Figure~\ref{fig:ch17_design_space}
combines the first three constraints into a single measure, the spectral efficiency
$\frac{T}{T+T_{\text{cp}}}\log_2(1+\text{SINR})$, where the SINR includes the Doppler ICI, for four
scenarios at an SNR of 25\,dB. Each curve has a broad optimum. For an urban macro-cell with a
4.7\,$\mu$s prefix and 200\,Hz Doppler the optimum is near 12\,kHz, and 15\,kHz costs almost nothing. For a
high-speed train at 1.5\,kHz Doppler it moves to about 40\,kHz, which is one reason NR allows 30\,kHz
spacing in FR1. For a broadcast single-frequency network with a 200\,$\mu$s prefix, the optimum is below
2\,kHz: DVB-T2 uses spacings of a few hundred hertz. Indoor Wi-Fi, with a short prefix and negligible
Doppler, has a flat optimum over two decades. The 312.5\,kHz of 802.11a, chosen to keep the FFT small
(64 points) and the symbol short for bursty packet traffic, sits well to the right of it and pays a
20\% prefix overhead. That is exactly the inefficiency 802.11ax removed by cutting the spacing by a
factor of four (Section~\ref{sec:ch17:systems}).

\mdcfig{ch17_design_space}{The subcarrier-spacing trade-off. (a) Prefix overhead (solid) rises with
spacing because the symbol shortens; the share of impairment due to Doppler ICI (dashed) falls with
spacing. (b) The resulting spectral efficiency at an SNR of 25\,dB, with the optimum marked for each
scenario. Dotted verticals mark 15\,kHz (LTE) and 312.5\,kHz (802.11a). The model ignores phase noise,
latency and FFT cost, which all favour wider spacing.}

\begin{worked}[The 802.11a numerology]
802.11a (1999, carried over unchanged into 802.11g at 2.4\,GHz in 2003) was designed for indoor channels
with RMS delay spreads up to about 100--200\,ns in a 20\,MHz channel.

\emph{Prefix.} A maximum excess delay of about 800\,ns gives $T_{\text{cp}}=0.8\,\mu$s, 16 samples at
20\,MS/s.

\emph{FFT and spacing.} A 64-point FFT at 20\,MS/s gives $\Delta f=312.5$\,kHz and $T=3.2\,\mu$s; the full
symbol is 4\,$\mu$s and the prefix overhead 20\%. A 256-point FFT would have cut the overhead to 6\%, but in
1999 a 64-point FFT was cheaper, and short symbols suit short packets and allow a fast turnaround between
transmit and receive.

\emph{Occupancy.} 52 subcarriers ($\pm1$ to $\pm26$) are used: 48 carry data and 4 (at $\pm7$ and $\pm21$) carry
BPSK pilots for phase tracking. The occupied band is $53\times312.5\,\text{kHz}\approx16.6$\,MHz, leaving 1.7\,MHz
guard bands on each side.

\emph{Rates.} The 48 data subcarriers carry $48\times b\times R$ information bits per 4\,$\mu$s symbol, with $b$
bits per subcarrier and code rate $R$. BPSK at rate 1/2 gives $48\times1\times\tfrac12/4\,\mu\text{s}=6$\,Mb/s;
64-QAM at rate 3/4 gives $48\times6\times\tfrac34/4\,\mu\text{s}=54$\,Mb/s, the familiar headline rate.

\emph{Doppler.} At 5.2\,GHz and walking speed (1.5\,m/s) $f_D\approx26$\,Hz, so $f_DT\approx8\times10^{-5}$: Doppler
ICI is negligible. The real constraint on the spacing is the crystal oscillator. With a $\pm20$\,ppm
tolerance on each end, the offset can reach $40\times10^{-6}\times5.2\,\text{GHz}\approx208$\,kHz, two-thirds of a
subcarrier, and the preamble must estimate and remove it before the FFT (Chapter~\ref{ch:sync}).
\end{worked}

\begin{worked}[The LTE numerology]
LTE (3GPP Release 8, 2008) was designed for macro-cells at 0.7--2.7\,GHz, with channel bandwidths of
1.4, 3, 5, 10, 15 and 20\,MHz and mobile speeds up to 350\,km/h.

\emph{Prefix.} Urban and suburban macro-cells have maximum excess delays up to about 5\,$\mu$s; the normal
prefix is 144 samples at 30.72\,MS/s, $4.69\,\mu$s (160 samples, $5.21\,\mu$s, for the first symbol of each
slot, to make seven symbols fill exactly 0.5\,ms).

\emph{Spacing.} At 350\,km/h and 2.6\,GHz, $f_D=vf_c/c\approx843$\,Hz. With $\Delta f=15$\,kHz, $f_DT=0.056$, and
the Doppler ICI is about $(\pi\times0.056)^2/6\approx5.2\times10^{-3}$, an SIR of 23\,dB: tolerable for the
QPSK and 16-QAM that a user at that speed would receive. Halving the spacing to 7.5\,kHz would halve the
overhead but quadruple the ICI.

\emph{FFT.} With 15\,kHz spacing, $f_s=N\times15$\,kHz. The 20\,MHz carrier uses $N=2048$ and
$f_s=30.72$\,MS/s, eight times the 3.84\,Mchip/s of UMTS, so one reference clock serves both. Its 100
resource blocks of 12 subcarriers occupy 18\,MHz, 90\% of the channel. The 10\,MHz carrier uses $N=1024$ at
15.36\,MS/s and 600 subcarriers; the 1.4\,MHz carrier uses $N=128$ at 1.92\,MS/s and 72 subcarriers.

\emph{Extended prefix.} For large cells and for multicast in single-frequency networks (MBSFN), an extended
prefix of 512 samples ($16.7\,\mu$s) reduces a slot to six symbols. A further MBSFN mode with 7.5\,kHz
spacing and a $33.3\,\mu$s prefix was specified for dedicated broadcast carriers.
\end{worked}

\subsection{The 5G NR numerologies}

LTE's single numerology could not stretch from sub-1\,GHz macro-cells to 28\,GHz and beyond. At millimetre
waves the Doppler for the same speed is ten times larger, oscillator phase noise is far stronger, and
narrow beams shorten the effective delay spread to tens or hundreds of nanoseconds. NR therefore defines a
family of \term{numerologies} indexed by $\mu$ (3GPP TS 38.211, Section~4.2),
\begin{equation}
\Delta f=2^\mu\times15\,\text{kHz},\qquad T_{\text{cp}}\approx\frac{4.69\,\mu\text{s}}{2^\mu},\qquad
\text{slot}=\frac{1\,\text{ms}}{2^\mu}=14\text{ symbols},
\end{equation}
listed in Table~\ref{tab:ch17:nr} and drawn in Figure~\ref{fig:ch17_numerology}. The prefix
\emph{fraction} is the same for every numerology, about 6.7\%, because the prefix and the symbol scale
together; what changes is the absolute time scale, and hence the balance between delay-spread tolerance on
one side and Doppler, phase-noise tolerance and latency on the other.

\begin{table}[htbp]\centering\small
\caption{NR numerologies with the normal cyclic prefix (3GPP TS 38.211 Table 4.2-1, Release 17). The prefix
of the first symbol in every 0.5\,ms is longer by $16\kappa T_c=0.52\,\mu$s. The typical-use column is a
simplification of the band-specific rules of TS 38.101.}
\label{tab:ch17:nr}
\begin{tabular}{@{}c r r r c r l@{}}
\toprule
$\mu$ & $\Delta f$ (kHz) & $T$ ($\mu$s) & $T_{\text{cp}}$ ($\mu$s) & Slots/ms & Slot ($\mu$s) & Typical use\\
\midrule
0 & 15 & 66.67 & 4.69 & 1 & 1000 & FR1 (below 7.125\,GHz), data and SSB\\
1 & 30 & 33.33 & 2.34 & 2 & 500 & FR1 mid-band, data and SSB\\
2 & 60 & 16.67 & 1.17 & 4 & 250 & FR1 and FR2 data; extended CP option\\
3 & 120 & 8.33 & 0.59 & 8 & 125 & FR2 (24--52.6\,GHz), data and SSB\\
4 & 240 & 4.17 & 0.29 & 16 & 62.5 & FR2 synchronisation blocks only\\
5 & 480 & 2.08 & 0.15 & 32 & 31.25 & 52.6--71\,GHz (Release 17)\\
6 & 960 & 1.04 & 0.07 & 64 & 15.63 & 52.6--71\,GHz (Release 17)\\
\bottomrule
\end{tabular}
\end{table}

\mdcfig{ch17_numerology}{Half a millisecond of each NR numerology from $\mu=0$ to $\mu=4$. Doubling the
subcarrier spacing halves the symbol, the prefix and the slot, so there are $2^\mu$ slots of 14 symbols per
millisecond. Symbol boundaries of different numerologies line up (every symbol of numerology $\mu$ spans
exactly two of numerology $\mu+1$), which lets a carrier mix numerologies in different bandwidth parts.}

All NR timing is expressed in the basic time unit $T_c=1/(480\,\text{kHz}\times4096)\approx0.509$\,ns. The
normal prefix is $144\kappa\,2^{-\mu}T_c$ with $\kappa=64$, and the useful symbol $2048\kappa\,2^{-\mu}T_c$; the
extra $16\kappa T_c$ on the first symbol of each half-subframe makes the symbols of all numerologies nest
exactly. A carrier may hold at most 275 resource blocks, 3300 subcarriers, so a 4096-point FFT suffices for
every configuration: 100\,MHz at 30\,kHz (273 resource blocks, 98.3\,MHz occupied, 122.88\,MS/s) or 400\,MHz
at 120\,kHz (264 resource blocks, 380\,MHz occupied, 491.52\,MS/s). The occupancy of up to about 98\% compares
with LTE's 90\%; it relies on the transmitter's spectral shaping (Section~\ref{sec:ch17:oob}), which the
standard leaves to the implementer as long as the emission masks of TS 38.101 and TS 38.104 are met.

\paragraph{Bandwidth parts.} A 100\,MHz NR carrier is too wide for a phone to monitor all the time, so NR
lets the network configure up to four \term{bandwidth parts} (BWPs) per carrier for each device, of which one
is active in each direction (TS 38.211 Section~4.4.5; TS 38.213). Each BWP is a contiguous set of resource
blocks with its own numerology and prefix. A phone can sit on a narrow 20\,MHz BWP while idle, saving the
power of a wide RF front end and a large FFT, and be switched by a downlink control message to the full
carrier when data arrives. Bandwidth parts also let a single carrier serve devices of different
capabilities, or different services with different numerologies, side by side.

\begin{inpractice}[Mixed numerologies and the guard between them]
Two numerologies on the same carrier are not orthogonal to each other: a 30\,kHz subcarrier is
not a whole number of cycles over a 15\,kHz symbol's FFT window, so energy leaks between the two
bandwidth parts. Implementations leave a small guard band (a few subcarriers) between them and apply
windowing or filtering per bandwidth part. This is one practical reason why the filtered-OFDM and WOLA
techniques of Section~\ref{sec:ch17:oob} were studied intensively during the design of NR, even though the
standard does not mandate any particular one.
\end{inpractice}

\subsection{Broadcasting: very long symbols}

Broadcasters face the opposite extreme. A digital TV network wants to reuse one frequency across a whole
country, with every transmitter sending the identical signal in a \term{single-frequency network} (SFN).
A receiver then sees the signals from several transmitters as echoes of one another, with relative delays set by
the path-length differences: tens of kilometres, or hundreds of microseconds. The prefix (which broadcasting
standards call the \emph{guard interval}) must cover them.

The 1995 DAB standard (ETSI EN 300 401) uses, in its Mode~I for terrestrial SFNs, 1536 subcarriers at 1\,kHz
spacing in 1.536\,MHz, a 1\,ms useful symbol and a 246\,$\mu$s guard interval: enough for transmitters about
70\,km apart. DVB-T (ETSI EN 300 744, 1997) offers a 2k and an 8k mode; in an 8\,MHz channel the 8k mode has
6817 subcarriers, a 896\,$\mu$s useful symbol, a spacing of about 1.12\,kHz and guard intervals of 1/32 to 1/4 of
the symbol (28 to 224\,$\mu$s). DVB-T2 (ETSI EN 302 755, 2008) extends this to FFT sizes up to 32k: in 8\,MHz, a
3584\,$\mu$s useful symbol, a spacing of about 279\,Hz and about 27\,000 subcarriers, with guard intervals up to
19/128 (532\,$\mu$s). ATSC~3.0 in North America (2017) uses similar 8k, 16k and 32k modes.

\begin{worked}[Sizing a DVB-T2 single-frequency network]
Transmitters of an SFN are spaced 60\,km apart. A receiver close to one transmitter sees the next one's
signal about 60\,km later: a relative delay of $60\,\text{km}/(0.3\,\text{km}/\mu\text{s})=200\,\mu$s. With the
32k mode in 8\,MHz ($T=3584\,\mu$s), a guard interval of 1/16 is $224\,\mu$s and covers it with some margin for
timing tolerance. The overhead is $1/17$, about 5.9\%. The same guard with the 8k mode ($T=896\,\mu$s) would be
a 1/4 guard interval and cost 20\%. The price of the long symbol is Doppler sensitivity: the spacing is only
279\,Hz, so a car at 120\,km/h receiving UHF channel 40 (626\,MHz), with $f_D\approx70$\,Hz, has
$f_DT\approx0.25$, far too much ICI for high-order QAM. The 32k mode is for rooftop reception; mobile
services use the smaller FFT sizes, and an SFN operator must choose between spacing its transmitters
widely and serving vehicles.
\end{worked}

\begin{history}[OFDM goes mainstream: 1987--2005]
The first mass-market OFDM systems were broadcast. In 1987 the European Eureka~147 project began
developing digital audio broadcasting; its design, based on work at the French CCETT laboratory in Rennes by
Michel Alard, Roselyne Lassalle and colleagues, used \emph{coded} OFDM (COFDM), in which a convolutional code
and time and frequency interleaving turn the channel's frequency selectivity into diversity
(Section~\ref{sec:ch17:coded}). DAB services began in the UK and Sweden in 1995, and DVB-T followed in 1997,
with the first services in 1998. On the telephone line, the 1993 Bellcore ``ADSL Olympics'' compared
single-carrier CAP modems with the discrete multitone (DMT) modem of Amati Communications, founded by
John Cioffi of Stanford; DMT won on performance and became the basis of ANSI T1.413 (1995) and ITU-T G.992.1
(1999). IEEE 802.11a adopted OFDM for 5\,GHz wireless LANs in 1999, with the same PHY as the European HiperLAN/2.
WiMAX (IEEE 802.16-2004 and 802.16e-2005) and Flarion's Flash-OFDM brought OFDMA to mobile broadband. When
3GPP studied the successor to UMTS in 2004--2005, OFDMA for the downlink was one of the less contested choices.
The uplink was debated longer, and in December 2005 3GPP chose single-carrier FDMA, DFT-spread OFDM, for its
lower peak-to-average ratio (Section~\ref{sec:ch17:dfts}). The first commercial LTE networks opened in
Stockholm and Oslo in December 2009.
\end{history}

%======================================================================
\section{The resource grid, pilots and channel estimation}
\label{sec:ch17:grid}
%======================================================================

\subsection{Thinking in time and frequency}

An OFDM signal is naturally described not as a waveform but as a two-dimensional \term{resource grid}:
subcarriers along one axis, OFDM symbols along the other. Each cell, one subcarrier for one symbol, is a
\term{resource element} (RE) and carries one complex number. LTE and NR group 12 consecutive subcarriers into a
\term{resource block} (RB), the unit in which the scheduler allocates spectrum, and group symbols into
slots of 14 (normal prefix). Figure~\ref{fig:ch17_grid} shows one slot of two resource blocks in NR. Each RE
is used for one of a small number of purposes: user data, control information, synchronisation signals,
or \emph{reference signals} (pilots) from which the receiver estimates the channel. The grid is where OFDM's
flexibility comes from. The scheduler can give different users different rectangles of it
(OFDMA, Section~\ref{sec:ch17:ofdma}), give different antennas different pilot patterns, and adapt the
pilot density to the channel.

\begin{figure}[htbp]\centering
\begin{tikzpicture}[x=0.5cm,y=0.19cm,font=\sffamily\scriptsize]
% control region
\fill[navy!12] (0,0) rectangle (2,24);
% data
\fill[practbg] (2,0) rectangle (14,24);
% DMRS type 1 on symbols 2 and 11, every other subcarrier
\foreach \s in {2,11}{\foreach \k in {0,2,...,22}{\fill[accent!75] (\s,\k) rectangle ++(1,1);}}
% one highlighted RE
\fill[navy] (6,17) rectangle ++(1,1);
\draw[step=1,gray!45,very thin,xstep=1,ystep=1] (0,0) grid (14,24);
\draw[navy,thick] (0,0) rectangle (14,24);
\draw[navy,thick] (0,12) -- (14,12);
% labels
\draw[decorate,decoration={brace,amplitude=4pt}] (-0.3,0) -- (-0.3,12) node[midway,left=5pt,align=right]{1 resource block\\= 12 subcarriers\\($12\Delta f$)};
\draw[decorate,decoration={brace,amplitude=4pt,mirror}] (0,-0.6) -- (14,-0.6) node[midway,below=5pt]{1 slot = 14 OFDM symbols ($1/2^\mu$\,ms)};
\draw[-{Stealth},navy] (7.3,17.5) -- (15.4,19.5) node[right,align=left]{one resource element:\\one subcarrier $\times$ one symbol};
\draw[-{Stealth},accent] (11.5,6.5) -- (15.4,8) node[right,align=left]{DMRS pilots (type 1,\\front-loaded + 1 additional)};
\draw[-{Stealth},navy!70] (1,22.5) -- (15.4,25) node[right,align=left]{control region (PDCCH)};
\node[right,align=left] at (15.4,13.5) {data (PDSCH)};
\draw[-{Stealth},practc] (15.35,13.5) -- (9,14.5);
\node[rotate=90] at (-6.2,12) {frequency};
\draw[-{Stealth}] (-5.6,4) -- (-5.6,20);
\end{tikzpicture}
\caption{The NR resource grid: one slot of two resource blocks, with a two-symbol control region, a
data region, and type-1 demodulation reference signals on every other subcarrier of the third and twelfth
symbols. A 100\,MHz carrier at 30\,kHz holds 273 such blocks, 3276 subcarriers, and 28 symbols per
millisecond: about 92\,000 resource elements every millisecond.}
\label{fig:ch17_grid}
\end{figure}

\subsection{Pilots as samples of the channel}

The receiver needs $H_k$ on every data RE. It measures it only where it knows what was sent, on the pilots,
and interpolates everywhere else. The channel $H(f,t)$ is a two-dimensional random function whose
variation is limited by the delay spread (in frequency) and the Doppler spread (in time)
(Chapter~\ref{ch:channels}). Pilot placement is therefore a sampling problem. By the sampling theorem,
pilots every $D_f$ subcarriers can represent channel impulse responses up to $1/(D_f\Delta f)$ long, and
pilots every $D_t$ symbols can represent Doppler spreads up to $\pm1/(2D_tT_s)$, where $T_s=T+T_{\text{cp}}$ is
the total symbol duration:
\begin{equation}
D_f\le\frac{1}{\tau_{\max}\,\Delta f},\qquad D_t\le\frac{1}{2f_D\,T_s}.
\label{eq:ch17:pilotnyq}
\end{equation}
In practice pilots are placed two to four times more densely than these limits, because interpolation
filters are not ideal and because noise on each pilot must be averaged out. The pilot overhead is
$1/(D_fD_t)$ of the grid for each antenna port.

Three families of pattern are in use (Figure~\ref{fig:ch17_pilots}).
\begin{description}[style=nextline,leftmargin=1.2em,font=\sffamily\bfseries\small\color{navy}]
\item[Block pilots.] Whole OFDM symbols of known data, sent as a preamble or at intervals. They sample the
channel densely in frequency and not at all in time, so they suit channels that are nearly static over a
packet: Wi-Fi's long training field, repeated as a \emph{midamble} in 802.11ax for long packets at speed.
\item[Comb pilots.] A few subcarriers carry pilots in every symbol. They track time variation well but
sample frequency sparsely. 802.11a/g/n/ac use four (or more, in wider channels) such pilots, not for channel
estimation but to track the common phase error from residual CFO and phase noise.
\item[Scattered pilots.] A diagonal lattice that shifts from symbol to symbol, so that a few pilots per
symbol sample both dimensions. DVB-T places a pilot on every twelfth subcarrier, shifted by three subcarriers
from symbol to symbol, so that four consecutive symbols together sample every third subcarrier, enough for echoes up to a third of the useful symbol,
more than the longest guard interval. LTE's cell-specific reference signals (CRS) for antenna port 0 occupy
every sixth subcarrier on two symbols per slot, staggered by three. NR replaces always-on reference signals
with user-specific ones sent only where data is scheduled.
\end{description}

\mdcfig{ch17_pilots}{Pilot patterns (red) on one slot of two resource blocks. (a) A block pilot symbol, as
in a Wi-Fi preamble. (b) Comb pilots on a fixed subset of subcarriers. (c) The scattered LTE CRS for one
antenna port. (d) NR type-1 DMRS on every other subcarrier of the third symbol, with one additional DMRS
symbol for mobility; the shaded symbols are the control region.}

\subsection{Least-squares estimation and interpolation}

On a pilot RE, $Y_p=H_pX_p+W_p$ with $X_p$ known, and the \term{least-squares} (LS) estimate is
\begin{equation}
\hat H_p^{\text{LS}}=\frac{Y_p}{X_p}=H_p+\frac{W_p}{X_p},\qquad \E|\hat H_p^{\text{LS}}-H_p|^2=\frac{N_0}{|X_p|^2}.
\end{equation}
It is unbiased and needs no statistics, but its error equals the inverse pilot SNR. Estimates between pilots
come from interpolation: linear interpolation in frequency and then in time is the simplest, and second-order
or spline interpolation is better on frequency-selective channels. Figure~\ref{fig:ch17_chest}a shows the
danger of linear interpolation near a deep fade, where the true $|H_k|$ has a narrow notch between pilots that
straight lines cannot follow; Figure~\ref{fig:ch17_chest}b shows the resulting error floor at high SNR.
Pilots are often transmitted with a few decibels of \emph{power boosting} to make their estimates cleaner,
paid for by slightly less power on the data.

\subsection{Exploiting the channel's structure}

The LS estimate ignores two things the receiver knows: the channel's impulse response is short (it fits
inside the prefix), and the channel's frequency response is therefore smooth. Two classic estimators use this.

\paragraph{DFT-based (transform-domain) estimation.} Transform the pilot LS estimates to the delay domain
with an IDFT, keep only the taps inside a window of length about $T_{\text{cp}}$ (where the channel must live),
set the rest to zero, and transform back to all subcarriers. The discarded taps contain only noise, so the
noise is reduced by roughly the ratio of the kept to the total taps: with 100 pilots spanning a window where
42 taps are kept, about 3.8\,dB. Two effects spoil the ideal. If the path delays are not integer multiples of the
delay resolution, each path's energy leaks into many taps, and truncation discards some of it. And if the
pilots do not span the whole FFT band, because of guard subcarriers, the leakage is worse at the band edges.
Both produce the error floor visible for the DFT-based estimator in Figure~\ref{fig:ch17_chest}b. The
estimator shines at low SNR and needs careful windowing at high SNR.

\paragraph{Linear MMSE estimation.} If the channel's frequency correlation is known, the best linear estimate
of the channel vector $\vect h$ on all subcarriers from the LS pilot estimates $\hat{\vect h}_p$ is the Wiener
filter
\begin{equation}
\hat{\vect h}^{\text{LMMSE}}=\vect R_{\vect h\vect h_p}\left(\vect R_{\vect h_p\vect h_p}+\frac{N_0}{E_p}\vect I\right)^{-1}\hat{\vect h}_p^{\text{LS}},
\label{eq:ch17:lmmse}
\end{equation}
where $\vect R_{\vect h\vect h_p}$ is the cross-correlation between the channel on all subcarriers and on the
pilots, and $\vect R_{\vect h_p\vect h_p}$ the pilots' autocorrelation. For a wide-sense-stationary uncorrelated
scattering channel with power delay profile $\{p_i,\tau_i\}$, the frequency correlation is
$\E[H_kH_l^*]=\sum_ip_ie^{-j2\pi(k-l)\Delta f\tau_i}$, the Fourier transform of the PDP (Chapter~\ref{ch:channels}).
The LMMSE estimator interpolates and denoises in one step. With the true PDP it gains 10--12\,dB over LS at
low SNR in Figure~\ref{fig:ch17_chest}b and has no floor.

The obvious objection is that the receiver does not know the PDP. Li, Cimini and Sollenberger (1998) showed
that a \emph{robust} design, which assumes a uniform PDP over the full prefix length and the maximum Doppler
the system supports, loses little compared with the matched estimator, because the estimator's performance
depends mainly on the support of the delay profile and Doppler spectrum, not their exact shape. The matrix
in~\eqref{eq:ch17:lmmse} can then be precomputed. Edfors and colleagues (1998) showed that a low-rank
approximation, via the eigen-decomposition of $\vect R_{\vect h\vect h}$, keeps nearly all of the gain at a
fraction of the cost, since only about $\tau_{\max}B$ eigenvalues are significant. Real receivers use such
separable, reduced-rank Wiener filters: one in frequency, one in time.

\mdcfig{ch17_chest}{Channel estimation on the EVA channel with a pilot on every sixth of 600 subcarriers.
(a) One realisation at a pilot SNR of 15\,dB (first 240 subcarriers): linear interpolation cuts through the
notch near subcarrier $-160$, while the LMMSE estimate follows it. (b) Mean-squared error versus pilot SNR.
Interpolation of LS estimates floors at high SNR; the DFT-based estimator gains at low SNR but floors
because the path delays are not on its delay grid; LMMSE with the true delay profile has neither
problem.}

\begin{inpractice}[Reference signals in NR]
NR has no always-on cell-wide reference signal like LTE's CRS, which occupied about 5--15\% of every
subframe (depending on the number of antenna ports) even when no one was being served and interfered with
neighbouring cells. Instead (3GPP TS 38.211, Section~7.4.1):
\begin{itemize}
\item The \emph{demodulation reference signal} (DMRS) is sent only in the resource blocks scheduled for a
user and is precoded like the data, so it measures the effective channel including beamforming. It is
\emph{front-loaded} (in the third or fourth symbol of the slot, or the first symbol of the allocation) so
that decoding can start early, and up to three \emph{additional} DMRS symbols can be configured for users at
high speed. Type~1 places pilots on every other subcarrier (a comb of two), with two code-division groups
that support up to four antenna ports with one symbol and eight with two; type~2 uses pairs of adjacent
subcarriers in three groups, for up to six or twelve ports.
\item The \emph{phase-tracking reference signal} (PT-RS) is a sparse comb, one subcarrier every two or four
resource blocks on every first, second or fourth symbol, configured at millimetre-wave frequencies to track the
common phase error of the oscillators (Section~\ref{sec:ch17:sync}).
\item The \emph{channel-state information reference signal} (CSI-RS) is used by the phone to measure the
channel for rank, precoder and modulation reporting (Chapter~\ref{ch:mimo}); the \emph{sounding reference
signal} (SRS) does the same on the uplink.
\end{itemize}
The design principle is ``transmit a pilot only when and where someone needs it,'' which reduces overhead,
interference and the base station's idle power consumption.
\end{inpractice}

\begin{worked}[Pilot density for a mobile channel]
A 3.5\,GHz NR carrier with 30\,kHz spacing ($T_s\approx35.7\,\mu$s) must serve users at 120\,km/h in a channel
with up to 3\,$\mu$s of excess delay. What pilot densities does~\eqref{eq:ch17:pilotnyq} require, and does
type-1 DMRS meet them?

\emph{Frequency:} $D_f\le1/(3\,\mu\text{s}\times30\,\text{kHz})=11$ subcarriers. Type-1 DMRS on every other
subcarrier ($D_f=2$) oversamples by more than five, leaving plenty of room for noise averaging, and also
supports multiple antenna ports by code division.

\emph{Time:} $f_D=33.3\,\text{m/s}\times3.5\,\text{GHz}/c\approx390$\,Hz, so $D_t\le1/(2\times390\times35.7\,\mu\text{s})
\approx36$ symbols. A single front-loaded DMRS per 14-symbol slot satisfies the sampling limit, but only
barely once the need for interpolation margin and the channel's variation within the slot are considered.
In practice, one additional DMRS symbol (Figure~\ref{fig:ch17_pilots}d) is configured above about 60--100\,km/h, and
two or three for high-speed trains. The overhead of one type-1 DMRS symbol is $\tfrac12\times\tfrac1{14}=3.6\%$; with
one additional symbol, 7.1\%.
\end{worked}

\begin{pitfall}[Estimating what you cannot see]
Channel estimates at the band edges are worse than in the middle, because interpolation there becomes
extrapolation. Pilots near guard bands, DC subcarriers and the edges of a user's allocation must be handled
specially, and DFT-based estimators suffer most. A related trap is to test a channel estimator only on
sample-spaced tapped-delay-line channels, whose paths fall exactly on the DFT delay grid. The estimator then
looks far better than it will on real channels, whose delays are continuous. The EVA and ETU models used in
this chapter round delays to the nearest sample at 15.36\,MS/s, but the DFT-based estimator of
Figure~\ref{fig:ch17_chest} works on a coarser delay grid set by the pilot spacing, which is why its floor appears.
\end{pitfall}

%======================================================================
\section{Synchronisation errors and what they do to OFDM}
\label{sec:ch17:sync}
%======================================================================

OFDM's orthogonality depends on the receiver's FFT window being aligned with the symbol, its
frequency grid with the subcarriers and its sample clock with the transmitter's. The algorithms that
achieve this---Schmidl--Cox preambles, cyclic-prefix correlation, the NR synchronisation signal
block---are covered in Chapter~\ref{ch:sync} (Sections~\ref{sec:ch10:freq} and~\ref{sec:ch10:frame}). This
section quantifies what happens when they are imperfect.

\subsection{Timing offset: the ISI-free window}

Suppose the FFT window starts $d$ samples \emph{before} its nominal position, inside the prefix. If
$0\le d\le N_{\text{cp}}-\nu$, the window still contains only samples of the current symbol, cyclically
shifted by $d$, and the circular-shift property of the DFT gives
\begin{equation}
Y_k=e^{-j2\pi kd/N}\,H_kX_k+W_k.
\end{equation}
The offset is harmless: it appears as a linear phase ramp across the subcarriers, which the channel
estimate absorbs as long as the pilots are processed with the same timing. There is therefore an
\emph{ISI-free window} of $N_{\text{cp}}-\nu+1$ acceptable starting positions (Figure~\ref{fig:ch17_timing_window}).
Starting later than nominal brings in samples of the next symbol; starting earlier than $N_{\text{cp}}-\nu$
brings in the previous symbol's tail. Either causes ISI and ICI whose power grows with the size of the
error. Because the penalty is asymmetric in practice (the first path is usually strong and the late
paths weak), receivers deliberately aim the window a few samples early, into the prefix. OFDM thus
needs only coarse timing, a striking contrast with single-carrier systems, which need the sampling phase
to a small fraction of a symbol (Chapter~\ref{ch:sync}).

\mdcfig[0.8]{ch17_timing_window}{Interference caused by the position of the FFT window (no noise; the $-60$\,dB
floor is the plotting limit). Anywhere in the green region the window contains a whole cyclic period of
the current symbol and the only effect is a phase ramp across subcarriers. Starting late or too early lets a
neighbouring symbol in.}

\subsection{Carrier frequency offset and inter-carrier interference}

A carrier frequency offset (CFO) of $\epsilon$ subcarrier spacings multiplies the received samples by
$e^{j2\pi\epsilon n/N}$. Substituting into the DFT,
\begin{equation}
\begin{split}
Y_k&=\sum_{l=0}^{N-1}H_lX_l\,S_{l-k}+W_k,\\
S_m&=\frac1N\sum_{n=0}^{N-1}e^{j2\pi(m+\epsilon)n/N}
=\frac{\sin\pi(m+\epsilon)}{N\sin\big(\pi(m+\epsilon)/N\big)}\,e^{j\pi(m+\epsilon)(N-1)/N}.
\end{split}
\label{eq:ch17:ici}
\end{equation}
The wanted term $S_0\approx\sinc(\epsilon)\,e^{j\pi\epsilon}$ is attenuated and rotated by the same amount
on every subcarrier; the terms $S_m$, $m\ne0$, leak each subcarrier into its neighbours
(Figure~\ref{fig:ch17_ici}a). By Parseval, $\sum_m|S_m|^2=1$, so the total ICI power, for independent unit-power
data on all subcarriers, is
\begin{equation}
P_{\text{ICI}}=1-|S_0|^2\approx1-\sinc^2\epsilon\approx\frac{(\pi\epsilon)^2}{3},\qquad
\text{SIR}\approx\frac{3}{(\pi\epsilon)^2}.
\label{eq:ch17:icipower}
\end{equation}
An offset of 1\% of the spacing limits the SIR to 35\,dB; 5\% to 21\,dB; 10\% to 15\,dB. Figure~\ref{fig:ch17_ici}b
compares~\eqref{eq:ch17:icipower} with simulation. Besides the ICI, the common rotation accumulates from
symbol to symbol by $2\pi\epsilon(N+N_{\text{cp}})/N$ radians, which the pilots must track.

The practical consequence is that OFDM receivers need a frequency estimate accurate to about 1--2\% of the
subcarrier spacing, which is 150--300\,Hz for LTE at 15\,kHz and only 0.1\,ppm at 2\,GHz. As
Chapter~\ref{ch:sync} explains, this is done in two stages: a coarse stage removes the \emph{integer} part of
$\epsilon$ (which shifts the whole grid by whole subcarriers without ICI) and the \emph{fractional} part
estimated from a repeated preamble or the prefix correlation; a tracking loop then follows the residual using
the pilots. On the OFDMA uplink the problem is harder, because every user's signal arrives with its own
offset, and the base station cannot correct them all with one rotation. Each device must therefore
pre-compensate its transmit frequency against the downlink, to within about $\pm0.1$\,ppm in LTE and NR.

\mdcfig{ch17_ici}{Inter-carrier interference. (a) Leakage coefficients $|S_l|^2$ from~\eqref{eq:ch17:ici} for
three frequency offsets ($N=64$): even a 5\% offset leaks $-25$\,dB into each adjacent subcarrier. (b) Signal-to-ICI
ratio: a frequency offset $\epsilon$ (theory~\eqref{eq:ch17:icipower}, circles), and a Rayleigh-fading channel with
Clarke's Doppler spectrum and maximum Doppler $f_D$ (theory~\eqref{eq:ch17:dopici}, squares). Doppler
spread is half as harmful as a frequency offset of the same size.}

\subsection{Doppler spread}

A time-varying channel produces ICI in the same way: the channel gain changes within the symbol, so
the subcarriers do not stay orthogonal. For a flat channel $h(t)$ that is constant in distribution with
autocorrelation $R_h(\Delta t)$, the useful term on subcarrier $k$ is the average of $h(t)$ over the
symbol and the rest is ICI. Its relative power is
\begin{equation}
P_{\text{ICI}}=1-\frac1{T^2}\int_0^T\!\!\int_0^TR_h(t-s)\,dt\,ds.
\end{equation}
For Clarke's model, $R_h(\Delta t)=J_0(2\pi f_D\Delta t)\approx1-(\pi f_D\Delta t)^2$ for small arguments
(Chapter~\ref{ch:channels}), and the average of $(t-s)^2$ over the square is $T^2/6$, so
\begin{equation}
P_{\text{ICI}}\approx\frac{(\pi f_DT)^2}{6},\qquad \text{SIR}\approx\frac{6}{(\pi f_DT)^2}.
\label{eq:ch17:dopici}
\end{equation}
This is half the ICI of a fixed offset $\epsilon=f_DT$, because the Doppler spectrum spreads the shifts
between $-f_D$ and $+f_D$ with a mean square of $f_D^2/2$. Li and Cimini (2001) give the tight upper bound
$(2\pi f_DT)^2/12$, twice~\eqref{eq:ch17:dopici}, valid for any Doppler spectrum. The approximation is the one used in
Figure~\ref{fig:ch17_design_space} and in the LTE worked example. For a high-speed train with a single
line-of-sight path, the Doppler is not a spread but a shift that changes sign as the train passes the mast;
it behaves as a CFO that the receiver can track, which is why NR's high-speed-train enhancements focus on
frequency tracking and pre-compensation rather than on subcarrier spacing.

\subsection{Phase noise: common phase error and ICI}

Oscillator phase noise $\phi(t)$ multiplies the signal by $e^{j\phi(t)}$ (Chapter~\ref{ch:sdr}). Expanding
it as a random, time-varying CFO, the DFT gives
\begin{equation}
Y_k=H_kX_k\,\underbrace{\frac1N\sum_ne^{j\phi[n]}}_{\text{CPE}}+\underbrace{\sum_{l\ne k}H_lX_l\,\frac1N\sum_n
e^{j\phi[n]}e^{j2\pi(l-k)n/N}}_{\text{ICI}}+W_k.
\end{equation}
The first term is a \term{common phase error} (CPE): the same rotation for every subcarrier of a symbol,
changing from symbol to symbol (Figure~\ref{fig:ch17_phase_noise}a). It is easy to remove by measuring the
rotation on pilots spread across the band, which is what 802.11's four pilots and NR's PT-RS are for. The
second term is ICI and cannot be removed so simply; it sets an EVM floor. For Wiener (free-running) phase
noise with a 3-dB linewidth $\beta$, the phase difference over $\Delta t$ has variance $2\pi\beta|\Delta t|$,
so $\E[e^{j(\phi(t)-\phi(s))}]=e^{-\pi\beta|t-s|}$. Repeating the Doppler calculation with this correlation,
and using $\E|t-s|=T/3$ over the square,
\begin{equation}
P_{\text{ICI}}\approx\frac{\pi\beta T}{3}=\frac{\pi}{3}\,\frac{\beta}{\Delta f}.
\end{equation}
The ICI power is proportional to the ratio of linewidth to spacing (Figure~\ref{fig:ch17_phase_noise}c). For
a practical locked oscillator the phase-noise spectrum is flat inside the loop bandwidth and falls outside
it, and the rule becomes: phase noise at offsets well below $\Delta f$ produces CPE, which pilots can remove,
while phase noise at offsets comparable to or beyond $\Delta f$ produces ICI, which they cannot. Since the phase noise of a frequency-multiplied oscillator rises by 6\,dB per doubling of its output
frequency, millimetre-wave NR uses wider
spacings (120\,kHz at 28\,GHz and up to 960\,kHz in the 60\,GHz band), which push more of the phase noise
into the correctable CPE region.

\mdcfig{ch17_phase_noise}{Wiener phase noise on 64-QAM OFDM (256-point FFT, linewidth 0.4\% of the spacing).
(a) Uncorrected, the common phase error rotates each symbol's constellation by a different angle, smearing it
into rings. (b) With the CPE removed, the remaining ICI blurs each point. (c) EVM versus linewidth: CPE
correction removes most of the error, leaving an ICI floor that follows $\pi\beta T/3$.}

\subsection{Sampling-clock offset and IQ imbalance}

If the receiver's ADC clock differs from the transmitter's DAC clock by a fraction $\delta$ (for example
$20\times10^{-6}$), two things happen. The FFT window drifts through the prefix by $\delta(N+N_{\text{cp}})$
samples per symbol, which must be corrected before it leaves the ISI-free window. And the whole
subcarrier grid is stretched: subcarrier $k$ appears with a frequency offset of $k\delta$ spacings, so it
rotates by $2\pi k\delta(N+N_{	ext{cp}})/N$ radians per symbol and suffers ICI that is strongest at the
band edges. For 802.11a at 20\,ppm, the drift is $20\times10^{-6}\times80=1.6\times10^{-3}$ samples per symbol, or one
sample after about 600 symbols (2.4\,ms): a long packet needs timing tracking. The phase ramp across the band
grows by about $2\pi\times26\times20\times10^{-6}\times80/64\approx0.004$ radians per symbol at the outermost
subcarrier; the pilots at $\pm7$ and $\pm21$ let the receiver estimate both the common phase and this slope.
In most receivers the carrier and sampling clocks derive from the same reference, so the two offsets are
estimated together.

\term{IQ imbalance} in a direct-conversion receiver (Chapter~\ref{ch:sdr}) produces a scaled copy of the
conjugate signal, and in OFDM terms this means that subcarrier $k$ is received with interference from its
mirror image, subcarrier $-k$: $Y_k=\mu H_kX_k+\nu H_{-k}^*X_{-k}^*+W_k$. An image rejection of 30\,dB limits
the subcarrier SNR to about 30\,dB, enough for 256-QAM but not for the 1024- and 4096-QAM of Wi-Fi~6 and~7. Such
receivers estimate and cancel the imbalance digitally, usually from the preamble.

\begin{worked}[A frequency-error budget at 28\,GHz]
An NR phone at 28\,GHz uses 120\,kHz spacing. Its crystal is disciplined to the base station's downlink with a
residual error of 0.1\,ppm, and it moves at 30\,km/h. What SIR do the offset and the Doppler allow, and how
would 15\,kHz spacing fare?

\emph{Offset:} $0.1\times10^{-6}\times28\times10^9=2.8$\,kHz, so $\epsilon=2.8/120=0.023$ and
$\text{SIR}\approx3/(\pi\times0.023)^2\approx570$, or 27.5\,dB.

\emph{Doppler:} $f_D=8.33\times28\times10^9/(3\times10^8)\approx780$\,Hz, $f_DT=0.0065$, and
$\text{SIR}\approx6/(\pi\times0.0065)^2$, about 41\,dB, negligible by comparison.

At 15\,kHz, the same offset would be 19\% of a spacing (SIR about 9\,dB) and the Doppler 5\% (SIR about 24\,dB):
unusable for anything above QPSK. The wide spacing is not a luxury at millimetre waves; it is what makes
the oscillators affordable.
\end{worked}

%======================================================================
\section{The peak-to-average power ratio}
\label{sec:ch17:papr}
%======================================================================

\subsection{Why OFDM is peaky}

Each time sample of an OFDM symbol is a sum of $N$ subcarrier contributions with independent data,
$x[n]=\frac1{\sqrt N}\sum_kX_ke^{j2\pi kn/N}$. For large $N$ the central limit theorem makes $x[n]$ nearly
complex Gaussian, so its envelope is Rayleigh distributed and its instantaneous power is exponentially
distributed: $\Pr\{|x[n]|^2>\gamma\bar P\}=e^{-\gamma}$. Most of the time the power is near its mean, but
occasionally many subcarriers add in phase and the signal peaks far above it
(Figure~\ref{fig:ch17_papr}a). The \term{peak-to-average power ratio} (PAPR) of a symbol is
\begin{equation}
\text{PAPR}=\frac{\max_n|x[n]|^2}{\E|x[n]|^2}\le N,
\end{equation}
with the upper bound reached only if all $N$ subcarriers align in phase: 18\,dB for $N=64$, 31\,dB for
$N=1200$. Such alignments are astronomically rare, and what matters is the distribution. If the $N$
Nyquist-rate samples are treated as independent, the complementary cumulative distribution function (CCDF) is
\begin{equation}
\Pr\{\text{PAPR}>\gamma\}\approx1-\big(1-e^{-\gamma}\big)^{N}.
\label{eq:ch17:paprccdf}
\end{equation}
The true analog peaks fall between the Nyquist samples, and van Nee and de Wild found empirically that
replacing $N$ by about $2.8N$ in~\eqref{eq:ch17:paprccdf} matches the oversampled signal well
(Figure~\ref{fig:ch17_papr}b, dotted). The key features are that the CCDF is steep, and that it moves to the right
only logarithmically with $N$: at a probability of $10^{-3}$ per symbol, the PAPR is about 10\,dB for 52
subcarriers, 10.8\,dB for 300 and 11.5\,dB for 1200. Every OFDM signal of practical size needs roughly 10--12\,dB
of headroom above its average power if its peaks are to pass undistorted.

\mdcfig{ch17_papr}{(a) The instantaneous power of three OFDM symbols with 200 QPSK subcarriers, relative to the mean;
the peak of each symbol is marked. (b) The PAPR CCDF, measured with 4$\times$ oversampling, for CP-OFDM with 52, 300
and 1200 subcarriers (the dotted lines are~\eqref{eq:ch17:paprccdf} with $2.8N$), for DFT-spread OFDM with 300
subcarriers, and for single-carrier QPSK with root-raised-cosine pulses ($\beta=0.22$).}

\subsection{Why it matters: the power amplifier}

A power amplifier (PA) is efficient only near saturation, and a signal with 10\,dB peaks must be backed off so
that its peaks stay in the PA's linear region (Chapter~\ref{ch:sdr}). Peaks that exceed it are compressed,
which has two effects: in-band distortion, measured as EVM, which limits the usable constellation; and
spectral regrowth outside the channel, which violates the emission mask and the adjacent-channel leakage
ratio (ACLR) limits. Of the two, spectral regrowth usually binds first, because the regulatory limits are
strict and the neighbour is someone else's customer. Peaks also cost resolution in the converters: every 6\,dB of
headroom costs one bit of DAC and ADC resolution.

\begin{worked}[What 10\,dB of back-off costs]
An ideal class-B amplifier has an efficiency of $\eta=\frac\pi4\sqrt{P_{\text{out}}/P_{\text{sat}}}$, 78.5\% at saturation.
Driving it with a signal whose average power is 10\,dB below saturation gives
$\eta=0.785\times10^{-10/20}=25\%$ for a sinusoid at that level. A class-A amplifier, whose efficiency is
$50\%\times P_{\text{out}}/P_{\text{sat}}$, would manage only 5\%. A 40\,W-average OFDM base-station PA backed
off by 10\,dB therefore needs 400\,W of peak capability, and at 25\% efficiency it draws 160\,W of DC power and
dissipates 120\,W as heat. Cutting the back-off to 6\,dB raises the class-B efficiency to 39\% and saves about
58\,W per amplifier. This is why every base-station transmitter combines \emph{crest-factor reduction}
(clipping and filtering, below) with digital predistortion (Chapter~\ref{ch:sdr}) and efficiency-enhancing
architectures such as the Doherty amplifier. For a phone, the same arithmetic appears as battery life and,
more importantly, uplink range: at the cell edge, the handset's transmit power is the limit, and every decibel
of back-off is a decibel of path loss the link cannot absorb.
\end{worked}

\subsection{Clipping and filtering}

The simplest way to reduce PAPR is to clip the signal's envelope at a level $A$, preserving the phase. Clipping
is a memoryless nonlinearity, and by Bussgang's theorem its output is a scaled copy of the input plus distortion
uncorrelated with it. The distortion falls both in band (raising the EVM) and out of band (splatter).
Filtering after clipping removes the out-of-band part, but it also regrows some of the peaks, so the
operation is repeated a few times (\term{iterative clipping and filtering}, Armstrong 2002). With
oversampling, the filter can be applied in the frequency domain simply by zeroing the unused bins.
Figure~\ref{fig:ch17_clipping} shows the trade-off for 16-QAM with 200 subcarriers clipped 4\,dB above the RMS
level. One pass of clipping and filtering reduces the PAPR at $10^{-3}$ from about 11\,dB to about 7.3\,dB with an
EVM of $-22$\,dB; four iterations reach about 5.7\,dB, at an EVM of $-19$\,dB. Without the filter, clipping splatters
energy 25\,dB below the in-band level right next to the channel, far above any emission mask.

\mdcfig{ch17_clipping}{Iterative clipping and filtering with the clipping level 4\,dB above RMS. (a) The PAPR
CCDF falls from about 11\,dB to 7.3\,dB after one clip-and-filter pass and to 5.7\,dB after four, at the cost of
in-band distortion (EVM in the legend). (b) Average per-symbol spectra: clipping alone splatters into the
neighbouring channels; the filter confines all versions to the occupied subcarriers.}

In a base station, \term{crest-factor reduction} (CFR) of this kind is applied to the composite multi-carrier
signal before the DPD, typically limiting the PAPR to 7--8\,dB. It is permitted because the standards budget
EVM for it: 3GPP TS 38.104 allows, for example, 8\% EVM for 64-QAM and 3.5\% for 256-QAM at the base-station
transmitter output, and the CFR consumes much of that budget. Smarter variants shape the clipping noise
into subcarriers that can tolerate it, such as unused or low-order-modulated ones.

\subsection{Distortionless methods: tone reservation, SLM and PTS}

A family of techniques reduces PAPR without distortion, by exploiting redundancy.
\begin{description}[style=nextline,leftmargin=1.2em,font=\sffamily\bfseries\small\color{navy}]
\item[Tone reservation (TR).] A few subcarriers (about 1--5\%) carry no data. The transmitter fills them with
values chosen, by a few iterations of a gradient or clipping-based algorithm, to cancel the peaks of the
time-domain signal. The receiver simply ignores them. Tellado and Cioffi developed it for DSL in the late 1990s,
and DVB-T2 includes it as an option, reserving about 1\% of the subcarriers.
\item[Active constellation extension (ACE).] The outer points of a QAM constellation can be moved
\emph{outwards}, away from the decision boundaries, without increasing the error probability. ACE (Krongold and
Jones, 2003) uses this freedom to reduce peaks. It costs a little average power and no data rate, and it is the
other PAPR option in DVB-T2.
\item[Selected mapping (SLM) and partial transmit sequences (PTS).] The transmitter generates several versions
of the same symbol, multiplying the data by different known phase sequences (SLM), or rotating blocks of
subcarriers by different phases (PTS), and sends the one with the lowest PAPR. With $U$ independent candidates,
the CCDF becomes $\Pr\{\text{PAPR}>\gamma\}^U$, a large improvement for small $U$. The receiver must learn which
candidate was sent, from side information or by blind detection, so these methods are popular in the
literature but rare in standards.
\end{description}
All of these reduce the PAPR at $10^{-3}$ by roughly 2--4\,dB. A more radical option is to change the waveform.

\subsection{DFT-spread OFDM: single carrier in OFDM clothing}
\label{sec:ch17:dfts}

\begin{figure}[htbp]\centering
\begin{tikzpicture}[node distance=0.34cm and 0.34cm,
  blk/.style={draw=navy,fill=navy!6,thick,minimum height=0.9cm,minimum width=1.15cm,align=center,font=\sffamily\scriptsize},
  key/.style={blk,fill=accent!8,draw=accent},
  arr/.style={-{Stealth[length=1.8mm]},thick,navy}]
\node[blk] (s) {$M$ QAM\\symbols};
\node[key,right=of s] (dft) {$M$-point\\DFT};
\node[blk,right=of dft,minimum height=1.9cm] (map) {Map to $M$\\contiguous\\subcarriers\\(zeros\\elsewhere)};
\node[blk,right=of map] (ifft) {$N$-point\\IFFT};
\node[blk,right=of ifft] (cp) {Add\\cyclic\\prefix};
\node[blk,right=of cp] (rf) {DAC\\and PA};
\draw[arr] (s)--(dft); \draw[arr] (dft)--(map); \draw[arr] (map)--(ifft); \draw[arr] (ifft)--(cp); \draw[arr] (cp)--(rf);
\node[below=0.25cm of map,font=\sffamily\scriptsize,text=accent,align=center] {other users occupy\\other subcarriers};
\end{tikzpicture}
\caption{A DFT-spread OFDM (SC-FDMA) transmitter. Without the red DFT block it is an ordinary OFDM
transmitter; with it, the output is a single-carrier signal whose spectrum occupies the $M$ assigned
subcarriers. The receiver removes the prefix, takes an FFT, equalises each subcarrier with an MMSE weight,
and applies an $M$-point inverse DFT.}
\label{fig:ch17_dfts_tikz}
\end{figure}

\term{DFT-spread OFDM} (DFT-s-OFDM), known in LTE as \term{SC-FDMA}, precodes each block of $M$ data symbols
$d_0,\dots,d_{M-1}$ with an $M$-point DFT before mapping them onto $M$ contiguous subcarriers of an $N$-point
OFDM modulator (Figure~\ref{fig:ch17_dfts_tikz}). The DFT and IDFT almost cancel. If $N=QM$ and the $M$
subcarriers start at bin 0, the output samples are
\begin{equation}
\begin{split}
x[n]&=\frac{1}{\sqrt{NM}}\sum_{m=0}^{M-1}d_m\sum_{k=0}^{M-1}e^{j2\pi k(n/N-m/M)}\\
&=\frac{1}{\sqrt{Q}}\sum_{m=0}^{M-1}d_m\,\frac{\sin\big(\pi(n/Q-m)\big)}{M\sin\big(\pi(n/Q-m)/M\big)}\,e^{j\pi(M-1)(n/Q-m)/M}.
\end{split}
\end{equation}
At $n=Qm$ the kernel is 1 at $m$ and 0 at all other symbols, so $x[Qm]=d_m/\sqrt Q$: the transmitted
samples at every $Q$th instant are the data symbols themselves, and the samples in between are
interpolated by a periodic sinc (Dirichlet) kernel. The signal is a single-carrier QAM signal at symbol rate
$M\Delta f$, shaped by an ideal band-limiting pulse and shifted in frequency to wherever the $M$ subcarriers were
placed. Its envelope is that of single-carrier modulation, not a Gaussian sum, and its PAPR is several decibels
lower than OFDM's (Figure~\ref{fig:ch17_papr}b): about 8\,dB for QPSK and 9\,dB for 16-QAM at $10^{-3}$,
compared with about 11\,dB for OFDM. The remaining peaks come from the sinc pulse, which has no excess
bandwidth; a single-carrier signal with root-raised-cosine pulses of roll-off 0.22 is lower still, around
5\,dB.

The receiver is the SC-FDE receiver of Section~\ref{sec:ch12:fde}: FFT, one-tap MMSE equalisation on each of the
$M$ subcarriers, then an $M$-point IDFT that returns to the symbol domain. Because each data symbol is spread over
all $M$ subcarriers, a fade on a few subcarriers is averaged out rather than destroying particular symbols, but
the MMSE equaliser pays a noise-enhancement penalty for it, and DFT-s-OFDM is somewhat worse than OFDM with the
same code at high SNR on frequency-selective channels. Multiple users still share the band by frequency division,
each with its own set of subcarriers, hence ``SC-FDMA''. What DFT-s-OFDM gives up is the flexibility of
putting different users or pilots on arbitrary individual subcarriers, and, in LTE, the ability to transmit
non-contiguous allocations (later relaxed to two clusters in LTE-Advanced).

\begin{keyidea}[Why the uplink is different]
A base station's power amplifier is shared, mains-powered and cooled, and it transmits dozens of carriers whose
combined PAPR it must handle anyway: OFDM's high PAPR costs money but not coverage. A handset's amplifier is
battery-powered, and at the cell edge its maximum output power limits the link. Every decibel of back-off is a
decibel less path loss the uplink can tolerate. LTE therefore uses OFDMA on the downlink and SC-FDMA on the uplink.
NR (TS 38.211, Section~6.3.1.4) offers both on the uplink: CP-OFDM for high rates and MIMO near the base station,
and DFT-s-OFDM (``transform precoding'') for coverage, switchable by the network. For the most
coverage-limited links, NR adds $\pi/2$-BPSK, whose alternating rotations keep consecutive symbols 90$^\circ$
apart so that the envelope never passes near zero. Combined with mild frequency-domain spectral shaping, it
takes the PAPR a few decibels below that of DFT-s-OFDM with QPSK.
\end{keyidea}

\begin{inpractice}[How much PAPR reduction is worth]
The 3GPP evaluations for LTE compared waveforms by their \emph{cubic metric}, a better predictor of the
required PA back-off than PAPR, and found SC-FDMA with QPSK roughly 2--3\,dB better than OFDMA, with a
smaller advantage for 16-QAM. That translates directly into a few percent more cell radius for an
uplink-limited cell, or equivalently into fewer base stations for a given coverage target. It is a small number
that decided a major design choice. The same logic makes single-carrier waveforms dominant wherever the
amplifier is the constraint: satellite links (DVB-S2 and its successors), where the transponder runs near
saturation, and the 60\,GHz 802.11ad/ay single-carrier PHY.
\end{inpractice}

%======================================================================
\section{Spectral containment}
\label{sec:ch17:oob}
%======================================================================

\subsection{The sidelobe problem}

Each subcarrier of a CP-OFDM symbol is a sinusoid with a rectangular time window, and its spectrum is a
$\sinc$ whose sidelobes decay only as $1/f^2$ in power. Summed over hundreds of subcarriers, the spectrum
(Figure~\ref{fig:ch17_ofdm_symbol}b) has a flat top and slowly decaying skirts: 1\,MHz beyond the edge of an LTE
carrier, the rectangular OFDM spectrum is only about 30\,dB down, not enough for adjacent-channel limits,
such as the 45\,dB adjacent-channel leakage ratio required of an LTE base station (3GPP TS 36.104). This is why OFDM systems leave guard bands (LTE uses 90\% of the
nominal channel), and why every transmitter applies additional shaping. The standards specify the emission
limits, not the method, and the common methods are windowing and filtering.

\subsection{Windowing: WOLA}

The sidelobes decay slowly because the symbol turns on and off abruptly. Replacing the rectangular window by
one with raised-cosine ramps, overlapping the ramps of successive symbols, smooths the transitions. In
\term{weighted overlap-and-add} (WOLA), the transmitter extends each symbol by $W$ extra cyclic samples at the
end, multiplies the extended symbol by a window that ramps up over its first $W$ samples and down over its last $W$,
and overlaps the ramps of adjacent symbols. The ramps eat into the prefix, so the effective prefix is shortened by
about $W$, but the sidelobes now decay much faster. The receiver can apply a similar window before its FFT to
reject interference from neighbouring carriers. WOLA is cheap: a few multiplications per sample. With a ramp of
1.6\,$\mu$s (a third of the LTE prefix), Figure~\ref{fig:ch17_oob} shows the emission 1\,MHz outside the carrier falling from
about $-30$\,dB to below $-55$\,dB.

\subsection{Filtering: f-OFDM}

The alternative is to pass the OFDM signal through a sharp band-pass filter matched to the allocation.
\term{Filtered OFDM} (f-OFDM) uses a filter whose passband covers the allocated subcarriers and whose length is
comparable to or longer than the symbol, typically a windowed sinc. The filter's impulse response is longer than
the prefix, so it causes some ISI, but because the filter is flat across the allocation, the ISI it causes is small
and concentrated at the band edges. In Figure~\ref{fig:ch17_oob}, a 33\,$\mu$s filter (half a symbol) pushes the
emissions below $-90$\,dB within a few hundred kilohertz of the edge. The cost is computation (a long filter at the
full sample rate, or a polyphase implementation) and latency. f-OFDM allows guard bands of just a few subcarriers,
and that is what makes NR's 98\% occupancy possible.

\mdcfig[0.85]{ch17_oob}{Power spectral density of an LTE-10\,MHz-like signal (600 subcarriers of 15\,kHz) with a
rectangular window, with WOLA using a 1.6\,$\mu$s raised-cosine taper, and with f-OFDM using a Hann-windowed
sinc filter of 33\,$\mu$s. Dotted lines mark the edges of the 10\,MHz channel.}

\subsection{Alternatives to CP-OFDM and why 5G kept it}

The 5G study phase (2015--2016) examined a dozen proposed waveforms. The main contenders were the following.
\begin{description}[style=nextline,leftmargin=1.2em,font=\sffamily\bfseries\small\color{navy}]
\item[FBMC/OQAM.] \term{Filter-bank multicarrier} descends from Saltzberg's 1967 offset-QAM scheme. Each
subcarrier is shaped by a long, well-localised prototype filter (for example the PHYDYAS filter, spanning four symbols),
which gives very low sidelobes and needs no cyclic prefix. Orthogonality is maintained only in the real domain,
by staggering the real and imaginary parts by half a symbol. The price is that the \emph{intrinsic interference}
between neighbouring symbols is purely imaginary rather than zero, which complicates channel estimation and
many MIMO schemes (Alamouti coding among them), and the long filters lengthen short bursts.
\item[UFMC.] \term{Universal filtered multicarrier} filters groups of subcarriers (sub-bands of, say, one resource
block) with short filters, and uses a zero-padded guard instead of a prefix to absorb the filter tails.
\item[GFDM.] \term{Generalised frequency-division multiplexing} transmits a block of several symbols on each
subcarrier with a circular pulse shape, and uses one prefix per block of symbols rather than per symbol. It is
flexible but non-orthogonal in general, which puts interference cancellation in the receiver.
\end{description}

3GPP kept CP-OFDM (TR 38.802 records the study). The reasons are instructive. CP-OFDM is trivially compatible with
MIMO: each subcarrier is a flat MIMO channel. Its channel estimation is simple and well understood. It has the
lowest latency, since each symbol stands alone. The spectral-containment advantage of the alternatives largely
vanished once it was recognised that WOLA or f-OFDM could be applied to CP-OFDM \emph{transparently}, without
the receiver needing to know, meeting the same emission limits. The alternatives' remaining advantages, such as
robustness to asynchronous access and tolerance of timing offsets, did not outweigh a decade of accumulated
engineering around CP-OFDM.

\begin{history}[The 5G waveform debate]
Between about 2013 and 2016 the choice of the 5G waveform was one of the liveliest topics in communications
research. The European 5GNOW and PHYDYAS projects championed GFDM and FBMC; Alcatel-Lucent and Nokia promoted UFMC;
Huawei proposed filtered OFDM. Survey papers with titles such as ``Modulation formats and waveforms for 5G
networks: who will be the heir of OFDM?'' (Banelli and colleagues, \emph{IEEE Signal Processing Magazine}, 2014)
framed the question. The eventual answer, CP-OFDM with implementation-specific windowing and filtering, can look
like a victory for the status quo. It is better read as a finding that the weaknesses of OFDM could be fixed at the
transmitter without changing what the receiver sees, while its strengths, MIMO-friendliness and simple
equalisation, could not be matched. The study did change NR: the flexible numerology and the bandwidth-part
concept are, in part, answers to the problems the new waveforms were proposed to solve.
\end{history}

%======================================================================
\section{Coded OFDM, frequency diversity and bit loading}
\label{sec:ch17:coded}
%======================================================================

\subsection{Uncoded OFDM is a Rayleigh channel}

The identity $Y_k=H_kX_k+W_k$ has an uncomfortable consequence. On a Rayleigh-fading channel, each $H_k$ is a
complex Gaussian random variable, so each subcarrier is a flat Rayleigh channel, and the average bit error
rate of uncoded OFDM is the Rayleigh error rate of Chapter~\ref{ch:channels}: it falls only as $1/\text{SNR}$.
For QPSK, $P_b\approx1/(4\,E_b/N_0)$ at high SNR, so $10^{-4}$ needs about 34\,dB, roughly 26\,dB more than on an
AWGN channel. The subcarriers near the notches of the frequency response
(Figure~\ref{fig:ch17_one_tap}a) produce nearly all the errors. OFDM has removed the ISI but not the
fading, and a frequency-selective channel that a single-carrier equaliser would have exploited for diversity
has been turned into many independent flat fades.

\subsection{Coding across subcarriers}

The remedy is to spread each codeword across many subcarriers, so that a few faded subcarriers damage
only a few code bits, which the decoder corrects from the others. If the codeword spans $L$ independently
fading subcarriers and the code has free distance $d_{\text{free}}$, the diversity order, the slope of the
error-rate curve at high SNR, is approximately $\min(d_{\text{free}},L)$. The number of independent fades in
a bandwidth $B$ is roughly $B/B_c$, where $B_c$ is the coherence bandwidth. Coded OFDM therefore turns a wide
bandwidth and a large delay spread into diversity, recovering what the one-tap equaliser appeared to
throw away. This is the ``C'' in \term{COFDM}, the name the DAB and DVB-T designers used.

Two ingredients are essential. The first is an \emph{interleaver} that sends adjacent code bits to
subcarriers far apart in frequency, because adjacent subcarriers fade together and a convolutional
code, in particular, cannot correct a burst of errors. The second is \emph{channel-state information in the
decoder}: the soft demapper must compute the bit LLRs with the per-subcarrier gain $|H_k|^2$, so that bits
from faded subcarriers are treated as erasures rather than as confident mistakes. The combination, a
binary code, a bit interleaver and a soft demapper, is \term{bit-interleaved coded modulation} (BICM),
proposed by Ephraim Zehavi in 1992 and analysed by Giuseppe Caire, Giorgio Taricco and Ezio Biglieri in 1998. On
fading channels, BICM's diversity order is set by the Hamming distance of the binary code rather than the
Euclidean distance of the coded modulation, which is why it has displaced trellis-coded modulation in
every modern wireless system: 802.11 from 802.11a onward, DVB-T and T2, LTE and NR.

Figure~\ref{fig:ch17_coded} shows the effect with the $K=7$ rate-1/2 convolutional code of 802.11a
(Chapter~\ref{ch:classiccodes}) and QPSK on 600 subcarriers, each codeword filling one OFDM symbol. On the
ETU channel, whose coherence bandwidth of about 200\,kHz gives some 45 independent fades across 9\,MHz,
coding with interleaving reaches $10^{-4}$ at about 10\,dB, within 7\,dB of AWGN and more than 20\,dB better than
uncoded. Without the interleaver, the same code loses about 7\,dB because the code bits of each trellis
section sit on adjacent, correlated subcarriers. On the EPA channel, whose coherence bandwidth of about
4.6\,MHz provides only two or three independent fades across 9\,MHz, even the interleaved code has too little
diversity to exploit and needs about 17\,dB. The lesson is general: coded OFDM's performance depends on the
product of bandwidth and delay spread, and a narrow allocation in a low-dispersion channel must find its
diversity elsewhere, in time (interleaving across symbols, retransmission with HARQ) or in space
(Chapter~\ref{ch:mimo}).

\mdcfig[0.85]{ch17_coded}{Frequency diversity through coding. Bit error rate of QPSK on 600 subcarriers
with the $K=7$, rate-1/2 convolutional code and soft (max-log) LLRs that include the channel gains, one
codeword per OFDM symbol and an independent channel realisation per symbol. The interleaved code on ETU gains
more than 20\,dB over uncoded transmission at $10^{-3}$; removing the interleaver, or moving to the
less dispersive EPA channel, costs 6--7\,dB.}

\begin{inpractice}[The 802.11a interleaver and DVB-T2's rotated constellations]
802.11a interleaves the coded bits of each OFDM symbol with a two-stage block permutation (IEEE 802.11,
Clause 17). The first stage writes the bits into a matrix row by row and reads them out column by column, so that
adjacent coded bits land on subcarriers about three apart in the 48-subcarrier symbol; the second stage
alternates adjacent bits between the more and less significant positions of the QAM label, which have
different reliabilities. The deinterleaver undoes both before the Viterbi decoder. DVB-T2 adds a further
diversity trick, \emph{rotated constellations}: each QAM point is rotated so that its I and Q components each
carry information about both bits, and I and Q are then sent on different subcarriers. If one subcarrier fades
completely, the other component alone still identifies the point. This \emph{signal-space diversity} gives a
worthwhile gain for high-rate codes on channels with erasures, such as single-frequency networks with strong
echoes.
\end{inpractice}

\begin{pitfall}[Forgetting the channel state in the demapper]
A surprisingly common simulation mistake, and an occasional implementation one, is to equalise with
$Y_k/H_k$ and then compute LLRs as if the equalised symbol had the same noise variance on every subcarrier. The
noise on a faded subcarrier has been amplified by $1/|H_k|^2$; ignoring this gives confident LLRs for
garbage bits, and the decoder is misled. The loss on a frequency-selective channel is several decibels,
comparable to throwing away the soft information altogether. The correct LLRs follow from $Y_k=H_kX_k+W_k$
directly, or from the equalised symbol with noise variance $N_0/|H_k|^2$.
\end{pitfall}

\subsection{Bit loading and water-filling}

When the transmitter knows the channel, it can do better than spread every codeword uniformly: it can put more
bits on strong subcarriers and fewer, or none, on weak ones. This is \term{adaptive bit loading}, and it is
the defining feature of \term{discrete multitone} (DMT), the name OFDM goes by on telephone wires. A DSL
line is almost static, so the receiver measures the SNR of every subcarrier (``tone'') during training,
computes a bit allocation, and sends it back to the transmitter.

The optimum allocation of power follows from information theory (Chapter~\ref{ch:infotheory}). With $N$
parallel Gaussian channels of gain-to-noise ratios $g_k=|H_k|^2/N_{0,k}$ and a total power $P$, the rate
\begin{equation}
R=\sum_k\log_2\!\left(1+\frac{P_kg_k}{\Gamma}\right)\quad\text{subject to}\quad\sum_kP_k=P
\end{equation}
is maximised, by a Lagrange-multiplier argument, when
\begin{equation}
P_k=\left(\mu-\frac{\Gamma}{g_k}\right)^{\!+},
\label{eq:ch17:waterfill}
\end{equation}
where the ``water level'' $\mu$ is chosen to use the total power. This is \term{water-filling}: picture the
curve $\Gamma/g_k$ as the bottom of a vessel, pour in power, and each subcarrier gets the depth of water above it
(Figure~\ref{fig:ch17_waterfill}a). Subcarriers whose floor is above the water level get nothing. The
\term{SNR gap} $\Gamma$ converts capacity into a practical rate: uncoded QAM at a symbol error rate of $10^{-7}$
needs an SNR about 9.8\,dB above the Shannon limit, a code reduces the gap by its coding gain, and a system
margin (6\,dB is common in DSL) increases it. The number of bits on tone $k$ is then
$b_k=\log_2(1+\text{SNR}_k/\Gamma)$, rounded down to an integer.

In practice, water-filling's power allocation matters less than its choice of \emph{which} tones to use. DSL
transmitters are constrained by a PSD mask on every tone, so they cannot pour extra power into good tones, and
theory and experience agree that a flat PSD over the tones water-filling would switch on comes very close to
the optimum. The bits, not the power, are what adapt. Practical loading algorithms work in discrete bits:
Hughes-Hartogs' greedy algorithm (1987) adds one bit at a time to the tone where it costs least power; the
algorithm of Chow, Cioffi and Bingham (1995) is a fast approximation; and the Levin--Campello algorithm finds the
exact discrete optimum efficiently.

\mdcfig{ch17_waterfill}{Bit loading on an ADSL2+ loop (tones 33--511 at 4.3125\,kHz, flat $-40$\,dBm/Hz transmit
PSD, a simple loss model for 0.4\,mm cable, crosstalk that rises with frequency and a narrowband AM
radio interferer near 1.35\,MHz; gap 11.8\,dB including a 6\,dB margin and 4\,dB coding gain). (a) Water-filling
on a 3\,km loop: the water level and the power allocation (blue). Tones above about 1.1\,MHz are switched off.
(b) The integer bit allocation for three loop lengths and the resulting line rates, at 4000 DMT symbols per
second.}

\begin{worked}[The ADSL2+ rate on a 3\,km loop]
ADSL uses a real-valued DMT signal: the IFFT input is Hermitian-symmetric, so that the output is a real baseband
signal for the copper pair. ADSL2+ (ITU-T G.992.5) uses a 1024-point IFFT at 4.416\,MS/s, giving 512 tones of
4.3125\,kHz up to 2.208\,MHz, of which tones 33--511 carry downstream data (the lower tones are for POTS
splitting and upstream). The prefix is 64 samples, 14.5\,$\mu$s, and after one synchronisation symbol per 68
data symbols the data-symbol rate is exactly 4000 per second. Each tone carries up to 15 bits.

For the 3\,km loop of Figure~\ref{fig:ch17_waterfill}, the SNR falls from about 60\,dB on the lowest tones to below the
gap near 0.9\,MHz. With the 11.8\,dB gap, the loading starts at 15 bits per tone, falls by roughly one bit every
30\,kHz, and reaches zero near 0.9\,MHz, for a total of about 1410 bits per DMT symbol: $1410\times4000\approx5.6$\,Mb/s.
On a 1\,km loop nearly every tone is usable and the line rate reaches about 27\,Mb/s before framing overheads
(the ADSL2+ maximum is about 24\,Mb/s of payload); at 5\,km only the lowest 70 tones survive, for about 1\,Mb/s.
The strong rate-versus-reach dependence is the defining characteristic of DSL, and the reason operators pushed
fibre ever closer to the home (Chapter~\ref{ch:wireline}).
\end{worked}

\begin{inpractice}[Adaptive modulation in wireless OFDM]
A mobile channel changes too quickly for tone-by-tone loading, and feeding back a bit allocation for thousands of
subcarriers every millisecond would cost too much. LTE and NR adapt more coarsely. The phone reports a
\emph{channel quality indicator} (CQI) for the whole band or for sub-bands of several resource blocks, and the base
station chooses one modulation and coding scheme for each user's allocation and relies on the code to average over
the fading within it. The frequency-domain adaptation happens instead in the \emph{scheduler}, which gives each
user the resource blocks where its channel is currently strong (Section~\ref{sec:ch17:ofdma}). Wi-Fi uses a
single MCS per packet (per resource unit in 802.11ax) chosen by rate-adaptation algorithms such as Minstrel, which
probe different rates and keep the one with the highest measured throughput.
\end{inpractice}

%======================================================================
\section{OFDMA, MIMO-OFDM and single-frequency networks}
\label{sec:ch17:ofdma}
%======================================================================

\subsection{OFDMA: sharing the grid}

\term{Orthogonal frequency-division multiple access} (OFDMA) gives different users different subsets of the
resource grid. On the downlink this is natural: the base station builds one IFFT input with each user's data
on its own resource blocks. Each user demodulates the whole symbol and keeps its own subcarriers. The payoff
is \term{multiuser diversity}. Different users see independent fading, so at any moment each resource block
is strong for someone, and a scheduler that assigns blocks to users with good channels, while keeping the
allocation fair over time (the \emph{proportional-fair} rule), gains substantially over a round-robin
allocation. Chapter~\ref{ch:multipleaccess} develops scheduling and its trade-offs.

The uplink is harder. Each user's signal travels over its own channel, with its own delay and its own
oscillator. For the subcarriers of different users to remain orthogonal at the base station's FFT, all
signals must arrive within the prefix of one another, and all must be on frequency to a small fraction of
the spacing. LTE and NR therefore command each phone's transmit timing with a \emph{timing advance}, which
compensates the round-trip delay so that signals from near and far users arrive aligned. Each phone also
locks its transmit frequency to the downlink it receives, and per-user power control keeps the strong signals
from overwhelming the leakage floor seen by weak ones. The same requirements appear in 802.11ax's uplink
OFDMA, where a trigger frame from the access point tells each station exactly when to transmit, and each
station must pre-correct its frequency and timing from that frame.

\subsection{MIMO-OFDM}

OFDM and multiple antennas fit together naturally. With $N_t$ transmit and $N_r$ receive antennas and a
prefix longer than every antenna pair's channel, each subcarrier becomes a flat MIMO channel,
\begin{equation}
\vect y_k=\vect H_k\vect x_k+\vect w_k,\qquad \vect H_k\in\C^{N_r\times N_t},
\end{equation}
where $\vect H_k$ is the DFT, element by element, of the matrix impulse response. All the MIMO techniques of
Chapter~\ref{ch:mimo}---spatial multiplexing, beamforming, space--time coding and multi-user MIMO---can then be
applied independently on each subcarrier, or on groups of subcarriers across which $\vect H_k$ changes
little. A single-carrier MIMO system on a frequency-selective channel faces a far harder problem, a matrix
equaliser with memory. This compatibility, more than any other single property, is why OFDM is the
waveform of MIMO systems: 802.11n (2009) brought MIMO-OFDM to the mass market, LTE built it in from the start,
and NR's massive-MIMO base stations with 64 or more antennas compute a precoder for every few resource blocks.

\subsection{Single-frequency networks}

Section~\ref{sec:ch17:design} introduced the single-frequency network from the point of view of the guard
interval. The idea is worth stating in general. In an SFN, the transmitters are synchronised in time and frequency
(usually by GNSS) and send identical OFDM symbols. The receiver cannot tell the signals of different transmitters
apart from multipath echoes, and as long as they arrive within the prefix, OFDM combines them automatically:
the channel $H_k$ is simply the sum of the channels from all transmitters. Far from being interference, a second
transmitter adds power and diversity, filling in the shadows and fades of the first. A country-wide TV or radio
network can then use one frequency instead of the four to seven a conventional network needs to keep
co-channel transmitters apart. SFNs are used by DAB and DVB-T/T2 throughout Europe, by ATSC~3.0, and by LTE's
multicast service (MBSFN), whose extended prefix and 7.5\,kHz mode were designed for them; later 3GPP releases
added still narrower numerologies with longer prefixes for broadcasting from high towers.

\begin{pitfall}[The SFN self-interference zone]
An echo from a distant transmitter that arrives \emph{later} than the guard interval is not merely wasted: it
is interference, with the ISI and ICI of~\eqref{eq:ch17:cpsir}. SFN planners must therefore check not only
the coverage of each transmitter but the delay differences everywhere in the service area, including over the
sea and from mountain-top sites that can be received at great distances under anomalous propagation. The
solution may involve adding deliberate static delays at some transmitters to move the zone of excessive delay
difference to where nobody lives.
\end{pitfall}

%======================================================================
\section{OFDM across the industry}
\label{sec:ch17:systems}
%======================================================================

Table~\ref{tab:ch17:systems} collects the parameters of the OFDM systems discussed in this chapter. Two things
stand out. The subcarrier spacings span more than three orders of magnitude, from the 279\,Hz of DVB-T2 to the
960\,kHz of NR at 60\,GHz, each set by the delay spread and Doppler of its channel. And the constellations have
grown steadily denser, to 4096-QAM in Wi-Fi~7 and DOCSIS~3.1, as the channels became cleaner (cable, short
indoor links) and the radios better.

\begin{table}[htbp]\centering\footnotesize
\caption{Parameters of some OFDM systems (maximum or typical values; see the cited standards for details). ADSL2+ uses a real-valued (Hermitian) DMT signal.}
\label{tab:ch17:systems}
\begin{tabular}{@{}l c l l l l@{}}
\toprule
System & Year & Bandwidth & FFT & $\Delta f$ & Densest QAM\\
\midrule
DAB (EN 300 401, mode I) & 1995 & 1.536\,MHz & 2048 & 1\,kHz & DQPSK\\
DVB-T (EN 300 744) & 1997 & 6--8\,MHz & 2k, 8k & 1.12--4.46\,kHz & 64\\
ADSL2+ (G.992.5) & 2003 & 2.2\,MHz & 1024 & 4.3125\,kHz & 15 bits/tone\\
IEEE 802.11a/g & 1999 & 20\,MHz & 64 & 312.5\,kHz & 64\\
IEEE 802.11n/ac & 2009/13 & 20--160\,MHz & 64--512 & 312.5\,kHz & 256\\
IEEE 802.11ax (Wi-Fi 6) & 2021 & 20--160\,MHz & 256--2048 & 78.125\,kHz & 1024\\
IEEE 802.11be (Wi-Fi 7) & 2024 & 20--320\,MHz & 256--4096 & 78.125\,kHz & 4096\\
DVB-T2 (EN 302 755) & 2008 & 1.7--10\,MHz & 1k--32k & down to 279\,Hz & 256\\
LTE (TS 36.211) & 2009 & 1.4--20\,MHz & 128--2048 & 15\,kHz & 256 (1024 later)\\
DOCSIS 3.1 downstream & 2013 & 24--192\,MHz & 4k, 8k & 50, 25\,kHz & 4096\\
5G NR (TS 38.211) & 2019 & 5--400\,MHz & up to 4096 & 15--960\,kHz & 256 (1024 in FR1)\\
\bottomrule
\end{tabular}
\end{table}

\subsection{Wi-Fi 6 and 7}

IEEE 802.11ax (Wi-Fi~6, approved in 2021 after products appeared in 2019) made three changes to the OFDM layer of
Wi-Fi. First, it quartered the subcarrier spacing to 78.125\,kHz with a 256-point FFT per 20\,MHz, making the
useful symbol 12.8\,$\mu$s. With a 0.8\,$\mu$s guard interval the overhead falls from 20\% to 5.9\%; with the 1.6 and
3.2\,$\mu$s options it can serve outdoor and large-venue channels. Second, it introduced OFDMA: the access point
can divide a channel into \term{resource units} (RUs) of 26, 52, 106, 242, 484, 996 or $2\times996$ subcarriers
and serve up to nine stations in one 20\,MHz transmission, which greatly reduces the overhead of the many short
packets typical of real networks. Third, it added 1024-QAM, which needs a transmitter EVM of about $-35$\,dB.
IEEE 802.11be (Wi-Fi~7, 2024) keeps the numerology, doubles the maximum channel to 320\,MHz, adds 4096-QAM
(EVM of about $-38$\,dB), allows several RUs to be assigned to one station (multi-RU), and allows a wide channel to
be used with holes in it (\emph{preamble puncturing}) when part of it is occupied by another network, a
flexibility that falls naturally out of OFDM, since a hole is just a set of zeroed subcarriers.

\begin{worked}[Wi-Fi 6 and 7 peak rates]
In 20\,MHz, 802.11ax uses 242 subcarriers: 234 for data and 8 for pilots. At the top MCS (1024-QAM, 10 bits,
rate 5/6) and a 13.6\,$\mu$s symbol (0.8\,$\mu$s guard), one spatial stream carries
$234\times10\times\tfrac56/13.6\,\mu\text{s}\approx143$\,Mb/s. A 160\,MHz channel has 1960 data subcarriers, so
$1960\times10\times\tfrac56/13.6\,\mu\text{s}\approx1.2$\,Gb/s per stream, and with the maximum of eight streams about
9.6\,Gb/s, the figure on Wi-Fi~6 marketing. For 802.11be at 320\,MHz, 3920 data subcarriers with 4096-QAM (12 bits,
rate 5/6) give about 2.9\,Gb/s per stream, and the standard's maximum of eight streams about 23\,Gb/s (the
46\,Gb/s often quoted came from early 16-stream proposals that did not survive into the final amendment). Real devices, with two to four antennas and
imperfect channels, see a small fraction of these numbers; but the arithmetic shows exactly which knobs the
standards turned: more subcarriers, more bits per subcarrier and more streams.
\end{worked}

\subsection{Cable, power line and the home network}

DOCSIS~3.1 (CableLabs, 2013) replaced the 6\,MHz single-carrier QAM channels of earlier cable modems with OFDM
channels up to 192\,MHz wide, using 4096- or 8192-point FFTs at 204.8\,MS/s, giving 50 or 25\,kHz subcarriers,
and a choice of prefix from 0.94 to 5\,$\mu$s to cover the echoes of the coaxial plant. The cable channel is
clean enough for 4096-QAM downstream, with 8192- and 16384-QAM as options, and the upstream uses OFDMA. Like DSL,
it adapts: the \emph{modulation profiles} assign different constellations to different groups of subcarriers
according to the measured SNR, a coarse form of bit loading. OFDM's ability to \emph{exclude} subcarriers is
equally valuable on the cable: subcarriers that coincide with strong ingress or with legacy channels are simply
switched off.

Power lines are a hostile medium: strong frequency-selective attenuation from impedance mismatches at every
outlet, impulsive noise from switching power supplies, and strict limits on radiation. Broadband power-line
standards (HomePlug AV and IEEE 1901, and ITU-T G.hn, G.9960) all use OFDM, with about 1100 subcarriers of
24.4\,kHz between 2 and 30\,MHz in HomePlug AV and bandwidths up to about 100\,MHz in later versions, with bit
loading per subcarrier and tone masks that notch out the amateur radio bands, a spectral agility no
single-carrier system could offer. Narrowband power-line standards for smart meters (G3-PLC and PRIME) use OFDM
too, at a few tens of kilohertz.

\subsection{Optical OFDM and Li-Fi}

Most short-reach optical links use intensity modulation and direct detection (IM/DD): the transmitted signal
is the optical \emph{power}, which must be real and non-negative. A real OFDM signal is produced by imposing
Hermitian symmetry on the IFFT input ($X_{N-k}=X_k^*$, with $X_0$ and $X_{N/2}$ real or zero), exactly as in DSL.
Non-negativity is obtained either by adding a DC bias large enough that clipping at zero is rare
(\emph{DC-biased optical OFDM}), or by loading only the odd subcarriers, which makes the signal antisymmetric
so that clipping its negative half at zero leaves the odd subcarriers intact, at half the spectral efficiency
(\emph{asymmetrically clipped optical OFDM}). DMT was studied intensively for 100G and 400G data-centre links
around 2015, but the industry chose the simpler PAM-4 (Chapter~\ref{ch:baseband}). OFDM thrives, however, in
visible-light and infrared \emph{Li-Fi}: IEEE 802.11bb (2023) adapts the 802.11 OFDM PHYs to light, using LEDs
or laser diodes as transmitters. For long-haul coherent optics, coherent optical OFDM had a period of intense
research around 2008--2012 but lost to single-carrier dual-polarisation QAM with powerful DSP
(Chapter~\ref{ch:wireline}), in part because the fibre's nonlinearity punishes OFDM's high PAPR.

\subsection{OFDM radar and integrated sensing}

An OFDM signal is also a very good radar waveform, and this observation underlies much of the current interest
in \emph{integrated sensing and communication} (ISAC) for 6G. Suppose the OFDM transmitter also listens for
reflections. A target at range $R$ and radial velocity $v$ returns the signal with delay $\tau=2R/c$ and Doppler
$f_d=2v/\lambda$. Because the transmitter knows the data $X_{k,l}$ on every subcarrier $k$ and symbol $l$, it can
divide it out of the received grid:
\begin{equation}
D_{k,l}=\frac{Y_{k,l}}{X_{k,l}}=\sum_pa_p\,e^{-j2\pi k\Delta f\tau_p}\,e^{j2\pi f_{d,p}lT_s}+\text{noise}.
\end{equation}
Each target is a two-dimensional complex exponential: an IDFT across subcarriers gives the range profile and a
DFT across symbols gives the Doppler, so a 2-D FFT of the grid produces a range--Doppler map. The range
resolution is $c/(2N\Delta f)=c/(2B)$ and the maximum unambiguous range is $c/(2\Delta f)$; the prefix must cover
the round-trip delay for the processing to remain ISI-free, which limits the range to about $cT_{\text{cp}}/2$ unless
longer processing is used. A 100\,MHz NR carrier at 30\,kHz spacing thus resolves targets 1.5\,m apart, with a
prefix-limited range of about 350\,m and, over 140 symbols (5\,ms) at 3.5\,GHz, a velocity resolution of
$\lambda/(2\times5\,\text{ms})\approx8.6$\,m/s. Automotive radar research has explored OFDM waveforms since Sturm
and Wiesbeck's work around 2009--2011, and 3GPP began studying ISAC in Release 19.

A related idea goes further: instead of placing data on the time--frequency grid, place it on the
\emph{delay--Doppler} grid, where a doubly dispersive channel is sparse and nearly constant. This is
\term{orthogonal time--frequency space} (OTFS) modulation, proposed by Hadani and colleagues in 2017, which
can be implemented as a two-dimensional precoding around an OFDM modem, much as DFT-s-OFDM is a one-dimensional
precoding. OTFS and its relatives, together with OFDM-based ISAC, are among the candidate 6G technologies examined
in Chapter~\ref{ch:sixg}.

%======================================================================
\section{Summary}
\label{sec:ch17:summary}
%======================================================================

\begin{itemize}
\item OFDM divides a wideband channel into $N$ orthogonal subcarriers spaced by $1/T$, whose spectra overlap but
cancel at each other's centres. The IDFT generates them and the DFT separates them (Weinstein and Ebert, 1971), at
a cost of $\tfrac12\log_2N$ multiplies per symbol.
\item A cyclic prefix at least as long as the channel memory makes the channel matrix circulant, which the DFT
diagonalises for every channel: $Y_k=H_kX_k+W_k$. Equalisation is one complex multiply per subcarrier. The prefix
costs $T_{\text{cp}}/(T+T_{\text{cp}})$ of time and energy.
\item The spacing is a compromise: small enough that the prefix overhead is low, large enough that Doppler
($P_{\text{ICI}}\approx(\pi f_DT)^2/6$), frequency offset ($\approx(\pi\epsilon)^2/3$) and phase noise
($\approx\pi\beta T/3$) cause little ICI. 802.11a (312.5\,kHz, 0.8\,$\mu$s), LTE (15\,kHz, 4.7\,$\mu$s), NR ($15\times2^\mu$\,kHz),
DVB-T2 (down to 279\,Hz, hundreds of microseconds) are each a solution for its channel.
\item The resource grid carries data, control and pilots. Pilot spacing must sample $H(f,t)$ at better than the
Nyquist rate set by the delay and Doppler spreads. LS estimates plus interpolation are simple; DFT-based and
LMMSE estimators exploit the channel's structure; NR sends demodulation pilots only with data.
\item Timing errors within the ISI-free part of the prefix only rotate the subcarriers. Frequency errors must be
kept to about 1--2\% of the spacing. Phase noise causes a common phase error, tracked by pilots, and ICI, which
sets a floor.
\item OFDM's PAPR is about 10--12\,dB at $10^{-3}$, costing amplifier efficiency and coverage. Clipping and filtering,
tone reservation and ACE reduce it; DFT-spread OFDM (SC-FDMA) is a single-carrier signal in OFDM form and is used
on the LTE and NR uplinks.
\item Rectangular symbols have slowly decaying sidelobes; WOLA and filtered OFDM contain the spectrum
transparently, which is one reason 5G kept CP-OFDM over FBMC, UFMC and GFDM.
\item Uncoded OFDM on a fading channel is a Rayleigh channel. Coding with bit interleaving (BICM) and CSI-weighted
LLRs recovers diversity of order up to $B/B_c$. When the channel is known at the transmitter, as in DSL, bit
loading and water-filling adapt the rate to each tone.
\item OFDMA shares the grid among users, MIMO-OFDM applies MIMO per subcarrier, and SFNs turn multiple transmitters
into benign echoes. OFDM carries Wi-Fi, cellular, broadcasting, DSL, cable, power-line and Li-Fi links, and it is
becoming a radar waveform too.
\end{itemize}

\begin{labbox}[OFDM from the IFFT to the resource grid]
The notebook \lab{moderncomms-labs/labs/lab07\_ofdm.py} builds an OFDM link step by step with
\texttt{commlib.ofdm}. It shows the orthogonality of overlapping subcarriers; equalises an EVA channel with one
tap per subcarrier (Figure~\ref{fig:ch17_one_tap}) and lets you shrink the cyclic prefix interactively until ISI
appears (Figure~\ref{fig:ch17_cp_length}); measures the ICI caused by a frequency offset against
$3/(\pi\epsilon)^2$ (Figure~\ref{fig:ch17_ici}); acquires timing and frequency with a Schmidl--Cox preamble;
compares comb-pilot LS channel estimation at three pilot spacings; compares the PAPR of CP-OFDM, DFT-s-OFDM and
single carrier (Figure~\ref{fig:ch17_papr}); and tabulates the NR numerologies. Its exercises extend it to LMMSE
channel estimation, Doppler ICI and clipping-and-filtering. The figure script \lab{book/figscripts/ch17\_figs.py}
generates every data figure in this chapter, including the coded-OFDM, water-filling, phase-noise and
spectral-containment experiments, and is a good starting point for the simulation problems below.
\end{labbox}

\problems
\begin{enumerate}[label=\thechapter.\arabic*]
\item Show that the subcarriers~\eqref{eq:ch17:subcarrier} are orthogonal over any interval of length $T$, not just
$[0,T)$. Why does this matter for the timing tolerance of Section~\ref{sec:ch17:sync}?
\item An OFDM system must cover a 40\,MHz channel with 90\% occupancy, a maximum excess delay of 2\,$\mu$s and a maximum
Doppler of 500\,Hz, with prefix overhead below 10\% and Doppler SIR above 25\,dB. Find the range of acceptable subcarrier
spacings, choose one, and give the FFT size, sample rate, prefix length in samples and number of used subcarriers.
\item Write out $\vect H_c$ for $N=4$ and $h=[1,0.5]$. Verify numerically that $\vect F\vect H_c\vect F^H$ is diagonal
with entries $H_k=1+0.5e^{-j2\pi k/4}$. What happens if the prefix is removed?
\item Prove that the zero-padded OFDM receiver, which adds the $\nu$ samples after the FFT window to the first $\nu$
samples of the window, produces a circular convolution. Compare the received SNR of ZP-OFDM and CP-OFDM for the same
transmitted energy per symbol.
\item Using~\eqref{eq:ch17:cpsir}, estimate the signal-to-interference ratio for the ETU profile with the NR normal
prefix at 30\,kHz spacing (2.34\,$\mu$s) and at 60\,kHz (1.17\,$\mu$s). Which numerology would you use for an ETU-like
macro-cell at 3.5\,GHz?
\item Derive~\eqref{eq:ch17:ici} and show that $\sum_m|S_m|^2=1$. For $\epsilon=0.1$ and $N=64$, compute $|S_0|^2$ and
$|S_{\pm1}|^2$ in decibels, and check against Figure~\ref{fig:ch17_ici}a.
\item Derive the Doppler ICI~\eqref{eq:ch17:dopici} from the two-dimensional integral, and repeat the calculation
for a channel whose Doppler spectrum is a single line at $+f_D$ (pure frequency offset). Explain the factor of two.
\item An 802.11ax station at 5.5\,GHz has a residual frequency error of 1\,ppm and uses 1024-QAM, which needs an SINR of
about 35\,dB. Is the ICI from this offset acceptable at 78.125\,kHz spacing? At 312.5\,kHz?
\item A DVB-T receiver in 8k mode with an 8\,MHz channel receives two SFN transmitters with a relative delay of
180\,$\mu$s. Which guard intervals (1/4, 1/8, 1/16, 1/32) avoid self-interference? What is the maximum vehicle speed at
690\,MHz for a Doppler SIR of 20\,dB?
\item Show that the PAPR of an OFDM symbol cannot exceed $N$, and that the bound is reached by setting all $X_k$ equal.
Using~\eqref{eq:ch17:paprccdf}, find the PAPR exceeded with probability $10^{-4}$ for $N=256$ and $N=4096$, with and
without the $2.8N$ correction.
\item Show that DFT-s-OFDM with $N=QM$ and localised mapping produces output samples $x[Qm]=d_m/\sqrt Q$, and that with
interleaved mapping (every $Q$th subcarrier) the output is the data sequence repeated $Q$ times, with a phase ramp.
What is the PAPR of the interleaved version for QPSK?
\item Derive the water-filling solution~\eqref{eq:ch17:waterfill} with a Lagrange multiplier. For three tones with
$g=[10,3,1]$, $\Gamma=1$ and $P=2$, find the powers and the rate. Compare with equal power.
\item \emph{(Simulation)} Extend \texttt{ch17\_figs.py}'s \texttt{coded()} function to 16-QAM with (a) Gray-labelled
BICM and the 802.11a interleaver and (b) no interleaver. Plot BER on EVA and ETU, and estimate the diversity order from
the slope.
\item \emph{(Simulation)} Implement an LMMSE channel estimator that assumes a uniform delay profile of length equal to the
prefix (the robust estimator of Li, Cimini and Sollenberger) and compare its MSE on EVA with the estimator matched to
the true profile (Figure~\ref{fig:ch17_chest}b).
\item \emph{(Simulation)} Generate an OFDM signal with 1200 subcarriers of 15\,kHz, apply iterative clipping and filtering
at 5, 6 and 7\,dB, and pass the result through the Rapp amplifier model in \texttt{commlib.channel.rapp\_pa} with a range of
back-offs. For each, measure the ACLR and the EVM, and find the back-off at which the ACLR reaches 45\,dB with and
without CFR.
\end{enumerate}

\furtherreading
\begin{itemize}
\item R.~van Nee and R.~Prasad, \emph{OFDM for Wireless Multimedia Communications}, Artech House, 2000. The classic
engineering text on OFDM, written by one of the designers of 802.11a; still the best practical treatment of
synchronisation, PAPR and coded OFDM.
\item Y.~Li and G.~L.~St\"uber (eds.), \emph{Orthogonal Frequency Division Multiplexing for Wireless Communications},
Springer, 2006. Chapters by leading researchers on channel estimation, ICI, PAPR, MIMO-OFDM and OFDMA.
\item S.~B.~Weinstein, ``The history of orthogonal frequency-division multiplexing,'' \emph{IEEE Communications
Magazine}, vol.~47, no.~11, 2009. A short first-hand history from Kineplex and Chang to LTE.
\item S.~B.~Weinstein and P.~M.~Ebert, ``Data transmission by frequency-division multiplexing using the discrete
Fourier transform,'' \emph{IEEE Trans. Communication Technology}, vol.~19, no.~5, 1971; and R.~W.~Chang, ``Synthesis
of band-limited orthogonal signals for multichannel data transmission,'' \emph{Bell System Technical Journal},
vol.~45, 1966. The two founding papers.
\item J.~A.~C.~Bingham, ``Multicarrier modulation for data transmission: an idea whose time has come,'' \emph{IEEE
Communications Magazine}, vol.~28, no.~5, 1990. The paper that relaunched multicarrier modulation, with the DSL view
of bit loading.
\item E.~Dahlman, S.~Parkvall and J.~Sk\"old, \emph{5G NR: The Next Generation Wireless Access Technology}, 2nd ed.,
Academic Press, 2020. Numerologies, bandwidth parts, reference signals and the reasoning behind them, by
members of the 3GPP design team. 3GPP TS 38.211 is the primary source.
\item Y.~Li, L.~J.~Cimini and N.~R.~Sollenberger, ``Robust channel estimation for OFDM systems with rapid dispersive
fading channels,'' \emph{IEEE Trans. Communications}, vol.~46, no.~7, 1998; and O.~Edfors, M.~Sandell, J.-J.~van de
Beek, S.~K.~Wilson and P.~O.~B\"orjesson, ``OFDM channel estimation by singular value decomposition,'' \emph{IEEE Trans.
Communications}, vol.~46, no.~7, 1998. The foundations of practical OFDM channel estimation.
\item T.~Pollet, M.~Van Bladel and M.~Moeneclaey, ``BER sensitivity of OFDM systems to carrier frequency offset and
Wiener phase noise,'' \emph{IEEE Trans. Communications}, vol.~43, 1995; and Y.~Li and L.~J.~Cimini, ``Bounds on the
interchannel interference of OFDM in time-varying impairments,'' \emph{IEEE Trans. Communications}, vol.~49, no.~3,
2001. The ICI analysis of Section~\ref{sec:ch17:sync}.
\item S.~H.~Han and J.~H.~Lee, ``An overview of peak-to-average power ratio reduction techniques for multicarrier
transmission,'' \emph{IEEE Wireless Communications}, vol.~12, no.~2, 2005. A clear survey of clipping, coding, SLM,
PTS, tone reservation and ACE.
\item G.~Caire, G.~Taricco and E.~Biglieri, ``Bit-interleaved coded modulation,'' \emph{IEEE Trans. Information Theory},
vol.~44, no.~3, 1998. The theory of BICM and why it suits fading channels.
\item T.~Starr, J.~M.~Cioffi and P.~J.~Silverman, \emph{Understanding Digital Subscriber Line Technology}, Prentice
Hall, 1999. DMT, bit loading and the practice of ADSL, from the people who built it.
\item B.~Farhang-Boroujeny, ``OFDM versus filter bank multicarrier,'' \emph{IEEE Signal Processing Magazine}, vol.~28,
no.~3, 2011; and P.~Banelli \emph{et al.}, ``Modulation formats and waveforms for 5G networks: who will be the heir of
OFDM?'', \emph{IEEE Signal Processing Magazine}, vol.~31, no.~6, 2014. The case for the alternatives.
\item C.~Sturm and W.~Wiesbeck, ``Waveform design and signal processing aspects for fusion of wireless communications
and radar sensing,'' \emph{Proceedings of the IEEE}, vol.~99, no.~7, 2011. OFDM radar and the origin of much of today's
ISAC research.
\end{itemize}
